\documentclass[11pt]{article}
\usepackage[margin=1in]{geometry}
\usepackage{amsmath,amssymb,amsthm}
\usepackage{enumitem}
\usepackage[hidelinks]{hyperref}

\newtheorem{theorem}{Theorem}[section]
\newtheorem{proposition}[theorem]{Proposition}
\newtheorem{lemma}[theorem]{Lemma}
\newtheorem{corollary}[theorem]{Corollary}
\newtheorem{sublemma}{Lemma}[theorem]
\renewcommand{\thesublemma}{\thetheorem\alph{sublemma}}
\newtheorem{citedinput}[theorem]{Input}
\theoremstyle{definition}
\newtheorem{definition}[theorem]{Definition}
\newtheorem{construction}[theorem]{Construction}
\theoremstyle{remark}
\newtheorem{remark}[theorem]{Remark}

\newcommand{\R}{\mathbb R}
\newcommand{\C}{\mathbb C}
\newcommand{\N}{\mathbb N}
\newcommand{\Z}{\mathbb Z}
\newcommand{\dd}{\,\mathrm d}
\newcommand{\Wtri}{W_{\triangle}}
\newcommand{\lean}[1]{\texttt{#1}}

\title{The logarithmic energy of $n$ points on the unit sphere:\\
\large the route to Corollary~1.3}
\author{Scott Armstrong \and Sylvia Serfaty}
\date{}

\begin{document}
\sloppy
\maketitle

\begin{abstract}
We prove, unconditionally,
\[
C_{\mathrm{tri}}\,n-o(n)\le E_{\log}(n)-\bigl(\tfrac12-\log2\bigr)n^2+\tfrac n2\log n\le C_{\rm ring}\,n+o(n),
\]
with $C_{\mathrm{tri}}=\tfrac12\log4\pi+\Wtri/\pi$, where $\Wtri$ is the torus constant of Theorem~1.2 of the
preprint~\cite{Preprint}, and $C_{\rm ring}$ an explicit $\eta$-product constant. Given the closed form of $\Wtri$
(statement (CF1) of Input~\ref{in:CS}, proved in the closed-form note), $C_{\rm ring}=C_{\mathrm{tri}}$, so
$E_{\log}(n)=(\tfrac12-\log2)n^2-\tfrac n2\log n+C_{\mathrm{tri}}\,n+o(n)$. Given also the value of
$|\eta(\rho)|$ (statement (CF2), also proved in the closed-form note), this is Corollary~1.3 of the paper.
The argument has four parts: an exact splitting identity in stereographic charts;
separation of minimizers by isotropic averaging; a lower bound by localization, screening (Proposition~7.18 of
Serfaty's \emph{Lectures on Coulomb and Riesz gases}~\cite{Lectures}), reflection onto square tori and Theorem~1.2 of the
preprint; and an upper bound by explicit
latitude-ring configurations. For the rings we prove an exact energy identity, which reduces the
upper bound to elementary per-row estimates. Every external input is stated with its reference.
\end{abstract}

\section*{Introduction}
Corollary~1.3 of the preprint~\cite{Preprint} gives the second term of the minimal logarithmic energy of $n$ points on the
unit sphere $S^2$: the constant $C_{\rm BHS}$ of Brauchart, Hardin and Saff~\cite{BrauchartHardinSaff2012}, which settles the case
$d=2$ of their conjecture. B\'etermin and Sandier~\cite{BeterminSandier2018} proved that the linear coefficient exists and
that it equals $C_{\rm BHS}$ exactly when the density-one triangular lattice minimizes the planar renormalized energy; the preprint
deduces Corollary~1.3 from this equivalence and Theorem~1.1. This note gives a different route to the same
asymptotics, which is the one followed in the formalization (\lean{SandierSerfaty.sphereLogEnergy\_asymptotics}). It does not pass
through the planar renormalized energy of $E_\triangle$: the lower bound localizes the energy of a minimizing configuration
in stereographic charts, screens it on squares, reflects it onto square tori and applies Theorem~1.2 of the preprint, which
is a statement about square tori; the upper bound is given by explicit configurations on latitude rings, whose energy is computed by an
exact identity. The constants are identified by the closed form of $\Wtri$ proved in the closed-form note
(Theorems~2.1, 4.7 and 5.8 there), and the screening step is the one proved in the screening note (Section~7 there).

Sections~\ref{sec:setting} and~\ref{sec:identities} fix the normalizations and prove the exact identities (splitting in charts,
conformal change of the puncture energy, truncation identity). Section~\ref{sec:separation} proves the separation of minimizers,
Section~\ref{sec:UB} the upper bound, and Section~\ref{sec:LB} the lower bound. Section~\ref{sec:gaps} lists the inputs and
dependencies.

\tableofcontents

%======================================================================
\section{Conventions and dependencies}\label{sec:status}

\paragraph{Sources.}
``Paper'' is the preprint~\cite{Preprint}. ``The Lectures'' is Serfaty's \emph{Lectures on Coulomb and Riesz gases}~\cite{Lectures};
numbers such as (7.1.21), Proposition~7.18 or Appendix~A refer to the Lectures. ``BS'' is B\'etermin--Sandier~\cite{BeterminSandier2018}.
Theorem~1.2 of the paper is \lean{SandierSerfaty.mul\_triangularEnergy\_le\_torusEnergy}.

\paragraph{Kinds of statements.}
Every numbered statement is \emph{proved here}, except those labelled ``Input'' (cited, with
reference). Section~\ref{sec:gaps} lists all inputs and dependencies.

\paragraph{Remarks on the route.}
\begin{itemize}[leftmargin=*]
\item The upper bound does not use stationarity of the triangular lattice. The exact identity of
  Proposition~\ref{prop:ring-identity} reduces the in-band distortion to a first-order error
  $O(T\sqrt n)=o(n)$. Stationarity, which improves the rate, is recorded in
  Remark~\ref{rem:stationarity}.
\item The construction produces every $n$, so the insertion bound
  (Lemma~\ref{lem:insertion}) is proved but not needed.
\item The torus constant enters the upper bound only through $C_{\rm ring}$; matching it with the
  lower bound uses (CF1).
\end{itemize}

\paragraph{Conventions.}
\begin{itemize}[leftmargin=*]
\item $\log$ is natural and $g(x)=-\log|x|$ on $\C$.
\item $B(x,r)$ is the open disc.
\item $f_\alpha(x):=\max\{0,\log(\alpha/|x|)\}$ for $\alpha>0$, so
  $\int_\C f_\alpha=\pi\alpha^2/2$. This is $f_\eta=(g-g(\eta))_+$ of the Lectures.
\item $C,c$ are constants whose dependence is stated; ``absolute'' means independent of all
  parameters.
\end{itemize}

%======================================================================
\section{Setting and normalizations}\label{sec:setting}

\begin{definition}[Energy]\label{def:energy}
$S^2=\{y\in\R^3:\|y\|=1\}$. For $Y=(y_1,\dots,y_n)\in(S^2)^n$, $n\ge2$, set
\[
E(Y)=\sum_{i\neq j}-\log\|y_i-y_j\|\in(-\infty,+\infty],
\qquad E_{\log}(n)=\min_{(S^2)^n}E,
\]
the sum running over \emph{ordered} pairs (paper, Corollary~1.3). Set
\[
\mathfrak K(Y):=E(Y)-\Bigl(\tfrac12-\log2\Bigr)n^2+\tfrac n2\log n .
\]
\end{definition}

\begin{definition}[Charts]\label{def:charts}
\begin{itemize}[leftmargin=*]
\item \emph{Stereographic chart} (from the north pole $\mathbf n=(0,0,1)$):
  $T(x)=\bigl(2\,\mathrm{Re}\,x,\,2\,\mathrm{Im}\,x,\,|x|^2-1\bigr)/(1+|x|^2)$ for $x\in\C$.
  It is a bijection $\C\to S^2\setminus\{\mathbf n\}$ with $T(0)=\mathbf s=(0,0,-1)$.
\item \emph{Conformal factor, density, potential:}
  \[
  \lambda(x)=\frac{2}{1+|x|^2},\qquad
  \mu(x)=\frac{\lambda(x)^2}{4\pi}=\frac{1}{\pi(1+|x|^2)^2},\qquad
  V(x)=\log(1+|x|^2).
  \]
\item \emph{Inversion:} $\iota(x)=1/\bar x$ on $\C\setminus\{0\}$. With
  $J(y_1,y_2,y_3)=(y_1,y_2,-y_3)$ we have $T\circ\iota=J\circ T$, so $x\mapsto T(\iota(x))$ is
  the chart from the south pole up to the isometry $J$.
\item \emph{Mercator map:}
  $\mathcal M(s,t)=(\operatorname{sech}t\cos s,\operatorname{sech}t\sin s,\tanh t)
  =T(e^{t+is})$.
\end{itemize}
\end{definition}

\begin{definition}[Puncture energy]\label{def:puncture}
Let $G\subset\C$ be open and $P=\{p_1,\dots,p_k\}\subset G$ finite. A function $H$ is
\emph{admissible on $(G,P)$} if:
\begin{itemize}[leftmargin=*]
\item $H\in C^2(G\setminus P)$;
\item near each $p_i$, $H-g(\cdot-p_i)$ extends to a $C^1$ function;
\item $\int_{G\setminus\bigcup_iB(p_i,\varepsilon)}|\nabla H|^2<\infty$ for every
  $\varepsilon>0$.
\end{itemize}
Then
\[
W_G(H):=\lim_{\eta\downarrow0}\Bigl[\tfrac12\int_{G\setminus\bigcup_iB(p_i,\eta)}|\nabla H|^2
+\pi k\log\eta\Bigr].
\]
\end{definition}

\begin{lemma}[The puncture limit exists]\label{lem:puncture-exists}
For admissible $H$ the limit in Definition~\ref{def:puncture} exists in $\R$. Moreover, if
$G=G_1\cup G_2\cup Z$ with $G_1,G_2$ open and disjoint, $Z$ Lebesgue-null and $Z\cap P=\emptyset$,
then $W_G(H)=W_{G_1}(H)+W_{G_2}(H)$.
\end{lemma}
\begin{proof}
Write $H=g(\cdot-p_i)+R_i$ near $p_i$ with $R_i\in C^1$. On the annulus
$A=B(p_i,\varepsilon)\setminus B(p_i,\eta)$,
\[
|\nabla H|^2=|x-p_i|^{-2}+2\nabla g\cdot\nabla R_i+|\nabla R_i|^2 .
\]
The first term integrates to $2\pi\log(\varepsilon/\eta)$. The other two terms are
$O(|x-p_i|^{-1})$, hence integrable on $B(p_i,\varepsilon)$. So
$\tfrac12\int_A|\nabla H|^2+\pi\log\eta$ converges as $\eta\downarrow0$. Additivity holds because
for small $\eta$ each removed disc lies in $G_1$ or in $G_2$, and $Z$ is null.
\end{proof}

\begin{definition}[Torus constant]\label{def:wtri}
$\Wtri:=\mathtt{triangularEnergy}$, the puncture energy of the mean-zero potential of one charge on
$\C/\Lambda_\triangle$, where $\Lambda_\triangle$ has covolume $1$. It equals the
paper's $W(E_\triangle)$ (Lemma~8.2 of the comparison note). Set
\[
C_{\mathrm{tri}}:=\tfrac12\log(4\pi)+\Wtri/\pi .
\]
\end{definition}

\begin{citedinput}[paper Theorem 1.2: the lower bound on square tori]\label{in:R1}
Let $N\ge2$ and $\Gamma_N=\sqrt N\,\Z[i]$. For $a_1,\dots,a_N\in\C$ with pairwise distinct classes
mod $\Gamma_N$,
\[
\mathtt{torusEnergy}(\Gamma_N,a)\ge N\,\Wtri .
\]
Here $\mathtt{torusEnergy}$ is the puncture energy over a fundamental domain of the mean-zero
torus potential $\mathtt{torusPotential}$, with $-\Delta h=2\pi(\sum_i\delta_{a_i}-1)$
(\lean{torusEnergy}, \lean{IsTorusPotential}). The potential exists
and is unique (\lean{existsUnique\_isTorusPotential}).
\end{citedinput}

\begin{definition}[Ring constant]\label{def:ering}
Let $\rho=\tfrac12+i\tfrac{\sqrt3}2$. Then $q:=e^{2\pi i\rho}=-e^{-\sqrt3\pi}$ and
$\eta(\rho)=q^{1/24}\prod_{m\ge1}(1-q^m)$. For $x>0$ put
\[
e(x):=-\tfrac12\log x+\frac{x}{12}-2\sum_{m\ge1}\log\bigl|1-(-1)^me^{-mx}\bigr| ,
\qquad
C_{\rm ring}:=e(\sqrt3\pi)+\tfrac12\log2 .
\]
Since $|q^{1/24}|=e^{-\sqrt3\pi/24}$,
\[
C_{\rm ring}=-\tfrac12\log\tfrac{\sqrt3\pi}{2}-2\log|\eta(\rho)| .
\]
The function $e$ is real-analytic on $(0,\infty)$.
\end{definition}

\begin{citedinput}[the closed form of $\Wtri$ and the value of $|\eta(\rho)|$; proved in the closed-form note]\label{in:CS}
\begin{itemize}[leftmargin=*]
\item (CF1) $\Wtri=-\pi\log\bigl(2\pi\sqrt{\sqrt3/2}\,|\eta(\rho)|^2\bigr)$.
\item (CF2) $|\eta(\rho)|=3^{1/8}\Gamma(1/3)^{3/2}/(2\pi)$.
\end{itemize}
Proofs: (CF1) is Theorem~4.7 of the closed-form note, together with Proposition~3.1 there and $|\Delta(\rho)|=|\eta(\rho)|^{24}$;
(CF2) is Theorem~5.8 there. See also B\'etermin--Sandier~\cite[\S7.3]{BeterminSandier2018}, whose
closed form contains the factor $2\sqrt2\cdot\pi$.
\end{citedinput}

\begin{lemma}[Constants]\label{lem:constants}
Under (CF1), $C_{\mathrm{tri}}=C_{\rm ring}$. Under (CF2),
$C_{\rm ring}=C_{\rm BHS}=2\log2+\tfrac12\log\tfrac23+3\log\frac{\sqrt\pi}{\Gamma(1/3)}$
(paper, \texttt{eq:bhs-constant}).
\end{lemma}
\begin{proof}
\emph{First claim.} By (CF1),
$\Wtri/\pi=-\log2\pi-\tfrac12\log\tfrac{\sqrt3}2-2\log|\eta(\rho)|$. Adding
$\tfrac12\log4\pi=\log2+\tfrac12\log\pi$ gives
$-\tfrac12\log\pi-\tfrac12\log\tfrac{\sqrt3}{2}-2\log|\eta(\rho)|=C_{\rm ring}$.

\emph{Second claim.} By (CF2), both sides equal
$\tfrac52\log2-\tfrac12\log3+\tfrac32\log\pi-3\log\Gamma(\tfrac13)$.

\emph{Numerics.} All three constants equal $-0.05560530494339251850092339872\ldots$
\end{proof}

\begin{theorem}[Main theorem]\label{thm:main}
Let $n\to\infty$ through all integers, and write $\min\mathfrak K=E_{\log}(n)-(\tfrac12-\log2)n^2+\tfrac n2\log n$.
\begin{enumerate}[label=(\roman*)]
\item \emph{Unconditionally,}
  \[
  C_{\mathrm{tri}}\le\liminf_n\frac{\min\mathfrak K}{n}\le\limsup_n\frac{\min\mathfrak K}{n}\le C_{\rm ring}.
  \]
\item \emph{Under (CF1),}
  \[
  E_{\log}(n)=\Bigl(\tfrac12-\log2\Bigr)n^2-\tfrac n2\log n+C_{\mathrm{tri}}\,n+o(n).
  \]
\item \emph{Under (CF1) and (CF2),} paper Corollary~1.3 holds with $C_{\mathrm{BHS}}$.
\end{enumerate}
\end{theorem}
\begin{proof}
\begin{itemize}[leftmargin=*]
\item Lower bound: Theorem~\ref{thm:LB} gives
  $\liminf_n\min\mathfrak K/n\ge C_{\mathrm{tri}}$.
\item Upper bound: Theorem~\ref{thm:UB} gives
  $\limsup_n\min\mathfrak K/n\le C_{\rm ring}$. This proves (i).
\item By Lemma~\ref{lem:constants}, $C_{\rm ring}=C_{\mathrm{tri}}$ under (CF1), which gives (ii);
  and $C_{\mathrm{tri}}=C_{\rm BHS}$ under (CF1) and (CF2), which gives (iii).\qedhere
\end{itemize}
\end{proof}

%======================================================================
\section{Exact identities}\label{sec:identities}

\begin{lemma}[Chordal identity]\label{lem:chordal}
For $x,x'\in\C$,
\[
\|T(x)-T(x')\|=\sqrt{\lambda(x)\lambda(x')}\,|x-x'| .
\]
For $x,x'\neq0$,
\[
|\iota(x)-\iota(x')|=\frac{|x-x'|}{|x|\,|x'|},
\qquad
\lambda(\iota(x))=|x|^2\lambda(x).
\]
For $(s,t),(s',t')\in\R^2$, writing $w=(t-t')+i(s-s')$,
\[
-\log\|\mathcal M(s,t)-\mathcal M(s',t')\|
=-\log|2\sinh(w/2)|+\tfrac12\log\cosh t+\tfrac12\log\cosh t' .
\]
\end{lemma}
\begin{proof}
\emph{Chordal identity.} Write $\|T(x)-T(x')\|^2=2-2\langle T(x),T(x')\rangle$ and expand; this
gives $4|x-x'|^2/((1+|x|^2)(1+|x'|^2))$.

\emph{Inversion.} Use $|1/\bar x-1/\bar x'|=|x'-x|/(|x||x'|)$ and
$\lambda(1/\bar x)=2|x|^2/(|x|^2+1)$.

\emph{Mercator.} With $x=e^{t+is}$:
\begin{itemize}[leftmargin=*]
\item $\lambda(x)=e^{-t}\operatorname{sech}t$;
\item $|e^{\zeta}-e^{\zeta'}|=e^{(t+t')/2}|2\sinh((\zeta-\zeta')/2)|$ for $\zeta=t+is$;
\item substitute into the chordal identity.
\end{itemize}
The identities were also checked numerically to $10^{-30}$, and against a brute-force
computation of the energy of row configurations.
\end{proof}

\begin{lemma}[Logarithmic potential of $\mu$]\label{lem:U}
$\int_\C\mu=1$. For all $x\in\C$,
\[
U(x):=\int_\C g(x-y)\mu(y)\dd y=-\tfrac12V(x).
\]
Moreover $\int V\dd\mu=1$ and $\int U\dd\mu=-\tfrac12$.
\end{lemma}
\begin{proof}
\emph{Circle averages.} The average of $g(x-re^{i\theta})$ over $\theta$ is
$-\log\max(|x|,r)$. With $R=|x|$ and $\mu$ radial,
\[
U(x)=-\log R\int_0^R\frac{2r\dd r}{(1+r^2)^2}-\int_R^\infty\frac{2r\log r}{(1+r^2)^2}\dd r .
\]

\emph{The two integrals.} The first is $R^2/(1+R^2)$. For the second, use
\[
\frac{2r\log r}{(1+r^2)^2}=\frac{\dd}{\dd r}\Bigl[-\frac{\log r}{1+r^2}
+\tfrac12\log\frac{r^2}{1+r^2}\Bigr];
\]
evaluating between $R$ and $\infty$ gives
$\frac{\log R}{1+R^2}-\log R+\frac12\log(1+R^2)$.

\emph{Conclusion.} Summing the two contributions gives $U(x)=-\tfrac12\log(1+R^2)$; the case
$R=0$ follows by continuity. Finally $\int V\dd\mu=\int_1^\infty s^{-2}\log s\dd s=1$, and
$\int U\dd\mu=-\tfrac12\int V\dd\mu=-\tfrac12$.
\end{proof}

\begin{definition}[Chart potential]\label{def:hX}
For pairwise distinct $X=(x_1,\dots,x_n)\in\C^n$ set
\[
h_X(x):=\sum_i g(x-x_i)+\tfrac n2V(x) .
\]
Then $-\Delta h_X=2\pi(\sum_i\delta_{x_i}-n\mu)$ on $\C$. As $|x|\to\infty$,
\[
h_X(x)=\sum_i\log\frac{|x|}{|x-x_i|}+\tfrac n2\log(1+|x|^{-2})=O(|x|^{-1}),
\qquad
\nabla h_X=O(|x|^{-2}).
\]
In particular $h_X$ is admissible on $(\C,X)$.
\end{definition}

\begin{proposition}[Exact splitting identity]\label{prop:splitting}
For pairwise distinct $X\in\C^n$,
\begin{align*}
W_\C(h_X)&=\pi\Bigl[\sum_{i\ne j}g(x_i-x_j)+n\sum_iV(x_i)-\tfrac{n^2}2\Bigr],\\
E(TX)&=\Bigl(\tfrac12-\log2\Bigr)n^2+\frac1\pi\Bigl[W_\C(h_X)+\pi\sum_i\log\lambda(x_i)\Bigr].
\end{align*}
\end{proposition}
\begin{proof}
\emph{Green's identity on the punctured disc.} Let $D_{R,\eta}=B(0,R)\setminus\bigcup_iB(x_i,\eta)$.
On $D_{R,\eta}$ we have $-\Delta h=-2\pi n\mu$. Green's identity gives
\[
\int_{D_{R,\eta}}|\nabla h|^2=-2\pi n\int_{D_{R,\eta}}h\mu+\int_{\partial B_R}h\,\partial_rh
-\sum_i\int_{\partial B(x_i,\eta)}h\,\partial_{r_i}h,
\]
where $r_i=|x-x_i|$.

\emph{Boundary terms.} The term at $R$ is $O(R^{-2})$. Write $h=g(\cdot-x_i)+H_i$ with $H_i$
smooth near $x_i$:
\[
H_i(x_i)=\sum_{j\ne i}g(x_i-x_j)-nU(x_i)=\sum_{j\ne i}g(x_i-x_j)+\tfrac n2V(x_i).
\]
Then $\int_{\partial B(x_i,\eta)}h\,\partial_{r_i}h=2\pi\log\eta-2\pi H_i(x_i)+o(1)$.

\emph{Limit.} Let $R\to\infty$ and $\eta\downarrow0$, and use
$\int h\mu=\sum_iU(x_i)-n\int U\dd\mu=-\tfrac12\sum_iV(x_i)+\tfrac n2$. This gives
\[
W_\C(h_X)=\pi\Bigl[\sum_iH_i(x_i)+\tfrac n2\sum_iV(x_i)-\tfrac{n^2}2\Bigr],
\]
which is the first identity.

\emph{Second identity.} By Lemma~\ref{lem:chordal}, $-\log\|Tx_i-Tx_j\|=g(x_i-x_j)-\tfrac12\log\lambda(x_i)-\tfrac12\log\lambda(x_j)$. Summing over ordered pairs,
\[
E(TX)=\sum_{i\ne j}g(x_i-x_j)-(n-1)\sum_i\log\lambda(x_i).
\]
Insert the first identity and $\log\lambda=\log2-V$:
\[
E(TX)=\frac{W_\C}{\pi}-\sum_iV(x_i)+\frac{n^2}{2}-n(n-1)\log 2
=\Bigl(\tfrac12-\log2\Bigr)n^2+\frac{W_\C}\pi+\sum_i\log\lambda(x_i).
\]

\emph{Checks.} The constants were also checked numerically (identity residual $10^{-16}$).
\end{proof}

\begin{lemma}[Scaling and conformal change of puncture energy]\label{lem:conformal}
Let $\Phi:\tilde G\to G$ be a $C^2$ diffeomorphism between open sets that is conformal or
anticonformal, with $|D\Phi(\tilde z)v|=k(\tilde z)|v|$. Let $H$ be admissible on $(G,P)$ and set
$\tilde P=\Phi^{-1}(P)$. Then $H\circ\Phi$ is admissible on $(\tilde G,\tilde P)$ up to an additive
constant near each point, and
\[
W_{\tilde G}(H\circ\Phi)=W_G(H)-\pi\sum_{\tilde p\in\tilde P}\log k(\tilde p).
\]
In particular, for $\Phi(z')=z'/\sigma$ ($\sigma>0$), $W_{\sigma G}(H(\cdot/\sigma))=W_G(H)+\pi|P|\log\sigma$.
\end{lemma}
\begin{proof}
\emph{Dirichlet invariance.} The Dirichlet integral is invariant under conformal and anticonformal
maps:
\[
\int_{\tilde G\setminus\Phi^{-1}(A)}|\nabla(H\circ\Phi)|^2=\int_{G\setminus A}|\nabla H|^2 .
\]

\emph{Images of small discs.} By Taylor expansion, for small $\eta$,
\[
B\bigl(p,k\eta(1-C\eta)\bigr)\subset\Phi(B(\tilde p,\eta))\subset B\bigl(p,k\eta(1+C\eta)\bigr),
\]
with $k=k(\tilde p)$ and $p=\Phi(\tilde p)$. On the annulus between these discs,
$|\nabla H|\le|z-p|^{-1}+C$, so its Dirichlet integral is at most
$\pi\log\frac{1+C\eta}{1-C\eta}+O(\eta)\to0$. Hence
\[
\tfrac12\int_{\tilde G\setminus\bigcup B(\tilde p,\eta)}|\nabla(H\circ\Phi)|^2
=\tfrac12\int_{G\setminus\bigcup B(p,k(\tilde p)\eta)}|\nabla H|^2+o(1).
\]
Adding $\pi|P|\log\eta=\pi\sum\log(k\eta)-\pi\sum\log k$ and letting $\eta\downarrow0$ gives the
claim.

\emph{Admissibility.} Near $\tilde p$, $g(\Phi(\tilde z)-p)=g(\tilde z-\tilde p)-\log k(\tilde p)+O(|\tilde z-\tilde p|)$, with a $C^1$ remainder.
\end{proof}

\begin{definition}[Blow-up]\label{def:blowup}
Fix the absolute constants
\[
c_{\rm sep}:=\sqrt{\log2/10},\qquad \kappa:=2/c_{\rm sep},\qquad s_n:=\kappa\sqrt n .
\]
For $X\in\C^n$ set $z_i=s_nx_i$ and $H(z):=h_X(z/s_n)$. Then
\[
-\Delta H=2\pi\Bigl(\sum_i\delta_{z_i}-m\Bigr),\qquad
m(z):=\frac{\lambda(z/s_n)^2}{4\pi\kappa^2},\qquad \int m=n .
\]
\end{definition}

\begin{corollary}[Blown-up identity]\label{cor:blown}
\[
\pi\,\mathfrak K(TX)=W_\C(H)+\pi\sum_i\log\lambda(x_i)-\pi n\log\kappa .
\]
\end{corollary}
\begin{proof}
By Lemma~\ref{lem:conformal}, $W_\C(H)=W_\C(h_X)+\pi n\log(\kappa\sqrt n)$. Insert this into
Proposition~\ref{prop:splitting}.
\end{proof}

\begin{lemma}[Inversion of the potential]\label{lem:inversion}
Let $X\in(\C\setminus\{0\})^n$ and $\tilde X=(\iota(x_i))_i$. Then on $\C\setminus\{0\}$,
\[
h_X\circ\iota=h_{\tilde X}-\sum_i\log|x_i| .
\]
\end{lemma}
\begin{proof}
By Lemma~\ref{lem:chordal},
$|\iota(\tilde x)-x_i|=|\tilde x-\tilde x_i|\,|x_i|/|\tilde x|$. Moreover
$V(\iota(\tilde x))=V(\tilde x)-2\log|\tilde x|$. Sum over $i$.
\end{proof}

\begin{proposition}[Two-chart decomposition]\label{prop:twochart}
Let $X\in\C^n$ be pairwise distinct with all $x_i\neq0$, and let $r>0$ with $|x_i|\ne r$ for all
$i$. Write $\mathcal S=\{i:|x_i|<r\}$ and $\mathcal N=\{i:|x_i|>r\}$, and use
Definition~\ref{def:blowup} for $X$ and for $\tilde X=\iota(X)$ (with $\tilde z_i=s_n\tilde x_i$,
$\tilde H(\tilde z)=h_{\tilde X}(\tilde z/s_n)$). Set $D=B(0,rs_n)$, $\tilde D=B(0,s_n/r)$, and
\[
\Phi_{\mathcal S}:=W_D(H)+\pi\sum_{i\in\mathcal S}\log\lambda(x_i),\qquad
\Phi_{\mathcal N}:=W_{\tilde D}(\tilde H)+\pi\sum_{i\in\mathcal N}\log\lambda(\tilde x_i).
\]
Then
\[
\pi\,\mathfrak K(TX)=\Phi_{\mathcal S}+\Phi_{\mathcal N}-\pi n\log\kappa .
\]
\end{proposition}
\begin{proof}
\emph{Splitting at $\partial D$.} By Lemma~\ref{lem:puncture-exists},
$W_\C(H)=W_D(H)+W_{\C\setminus\bar D}(H)$; no point lies on $\partial D$.

\emph{The outer part.} The map $\iota_n(\tilde z):=s_n^2/\bar{\tilde z}=s_n\iota(\tilde z/s_n)$
is an anticonformal diffeomorphism $\tilde D\setminus\{0\}\to\C\setminus\bar D$ with
$k(\tilde z)=s_n^2/|\tilde z|^2$. By Lemma~\ref{lem:inversion}, $H\circ\iota_n=\tilde H+\mathrm{const}$.
Since $0\notin\tilde X$, $\tilde H$ is smooth at $0$, so the puncture energy over
$\tilde D\setminus\{0\}$ equals that over $\tilde D$. Lemma~\ref{lem:conformal} then gives
\[
W_{\C\setminus\bar D}(H)=W_{\tilde D}(\tilde H)+\pi\sum_{i\in\mathcal N}\log\frac{s_n^2}{|\tilde z_i|^2}.
\]

\emph{Bookkeeping.} Since $|\tilde z_i|=s_n|\tilde x_i|$ and
$\lambda(x_i)=|\tilde x_i|^2\lambda(\tilde x_i)$ (Lemma~\ref{lem:chordal}), each
$i\in\mathcal N$ satisfies $\log\frac{s_n^2}{|\tilde z_i|^2}+\log\lambda(x_i)=\log\lambda(\tilde x_i)$.
Conclude with Corollary~\ref{cor:blown}.
\end{proof}

\begin{lemma}[Truncation identity]\label{lem:truncation}
Let $G\subset\C$ be open and $\varrho\in C^1(G)$. Let $H$ be admissible on $(G,P)$ with
$-\Delta H=2\pi(\sum_{p\in P}\delta_p-\varrho)$ in $\mathcal D'(G)$. Let $\alpha_p>0$ be such that
the closed discs $\bar B(p,\alpha_p)\subset G$ are pairwise disjoint, and set
$H_\alpha:=H-\sum_pf_{\alpha_p}(\cdot-p)$. Then
\[
W_G(H)=\tfrac12\int_G|\nabla H_\alpha|^2+\pi\sum_{p\in P}\log\alpha_p
-2\pi\sum_{p\in P}\int_{B(p,\alpha_p)}f_{\alpha_p}(x-p)\varrho(x)\dd x .
\]
\end{lemma}
\begin{proof}
\emph{Reduction to one disc.} Outside $\bigcup B(p,\alpha_p)$ we have $H_\alpha=H$. Fix $p$, write
$\alpha=\alpha_p$ and $f=f_\alpha(\cdot-p)$. In $B(p,\alpha)$, $f=g(\cdot-p)-g(\alpha)$, so
$H_\alpha\in C^1(B(p,\alpha))$ with $\Delta H_\alpha=2\pi\varrho$.

\emph{Expansion on the annulus.} On $A=B(p,\alpha)\setminus B(p,\eta)$,
\[
|\nabla H|^2=|\nabla H_\alpha|^2+2\nabla H_\alpha\!\cdot\!\nabla f+|\nabla f|^2,
\qquad
\tfrac12\int_A|\nabla f|^2=\pi\log(\alpha/\eta).
\]
For the cross term, integrate by parts. Since $f=0$ on $\partial B(p,\alpha)$, $f=\log(\alpha/\eta)$
on $\partial B(p,\eta)$, and $\Delta f=0$ in $A$,
\[
\int_A\nabla H_\alpha\!\cdot\!\nabla f
=-\log(\alpha/\eta)\int_{B(p,\eta)}\Delta H_\alpha-\int_Af\Delta H_\alpha .
\]
The first term is $O(\eta^2\log(1/\eta))$. The second term tends to
$-2\pi\int_{B(p,\alpha)}f\varrho$.

\emph{Limit.} Also $\int_{B(p,\eta)}|\nabla H_\alpha|^2=O(\eta^2)$. Adding $\pi\log\eta$ per point
and letting $\eta\downarrow0$ gives the claim.
\end{proof}

\begin{remark}[Comparison with the Lectures]\label{rem:book-F}
The truncated energies of the Lectures use the same normalization up to the factor $2\pi$:
$c_d=2\pi$, and (7.1.22)/(7.1.23) carry the factor $1/(2c_d)$. For
disjoint balls inside the domain, Lemma~\ref{lem:truncation} gives
\[
2\pi\,\mathsf F(\cdot)=W(\cdot).
\]
This is also the equality case of the monotonicity Lemma~4.13/(7.1.40).
\end{remark}

%======================================================================
\section{Minimizers: existence and separation}\label{sec:separation}

\begin{lemma}[Existence]\label{lem:existence}
$E$ is lower semicontinuous from $(S^2)^n$ to $(-\infty,+\infty]$ and is finite on configurations
of distinct points. Hence minimizers exist, and they consist of distinct points. Since $E$ is
invariant under the orthogonal group $O(3)$, a minimizer may be rotated so that no point lies at
two prescribed antipodal poles.
\end{lemma}
\begin{proof}
Each term $-\log\|y-y'\|\ge-\log2$ is lower semicontinuous with values in $(-\infty,\infty]$, and
$(S^2)^n$ is compact. For the last claim, the set of rotations $R$ with $R(Y)\cap\{\pm\mathbf n\}\ne\emptyset$
is a finite union of proper closed submanifolds of $SO(3)$.
\end{proof}

\begin{theorem}[Separation]\label{thm:separation}
Let $Y$ minimize $E$ on $(S^2)^n$, $n\ge2$. Then
\[
\|y_i-y_j\|\ge c_{\rm sep}/\sqrt n\quad\text{for all }i\neq j,
\qquad c_{\rm sep}=\sqrt{\log2/10}.
\]
Consequently, in any chart $x=T^{-1}(R\,y)$ ($R\in O(3)$, $R(Y)\not\ni\mathbf n$), we have
$|x_i-x_j|\ge c_{\rm sep}/(2\sqrt n)$. In blown-up coordinates (Definition~\ref{def:blowup}),
$|z_i-z_j|\ge1$.
\end{theorem}
\begin{proof}
\emph{Choice of chart.} Fix $i$. The closed hemisphere $\{p:\langle p,y_i\rangle\le0\}$ is
uncountable, so pick $p$ in it with $p\notin Y$, and a rotation $R$ with $Rp=\mathbf n$. Replacing
$Y$ by $RY$ (still a minimizer) and setting $x_k:=T^{-1}(y_k)$, we get $|x_i|\le1$.

\emph{Minimality in the chart.} By Lemma~\ref{lem:chordal},
\[
E(TX)=\sum_{a\neq b}g(x_a-x_b)+(n-1)\sum_aV(x_a)-n(n-1)\log2 .
\]
Its dependence on $x_i$ is $2\phi(x_i)$, where
\[
\phi(x)=\sum_{j\ne i}g(x-x_j)+\psi(x),\qquad \psi=\tfrac{n-1}2V .
\]
Every $x\in\C$ gives a point $T(x)\in S^2$, so minimality of $Y$ gives $\phi(x_i)\le\phi(x)$ for
all $x\in\C$.

\emph{Isotropic averaging.} Average over the circle $\partial B(x_i,r)$. The average of
$g(\cdot-x_j)$ is $-\log\max(r,d_j)$, where $d_j=|x_i-x_j|$. Since
$\Delta\psi=2(n-1)(1+|x|^2)^{-2}\le2(n-1)$, the mean-value formula gives
\[
\operatorname{avg}_{\partial B(x_i,r)}\psi-\psi(x_i)=\int_0^r\frac1{2\pi s}\int_{B(x_i,s)}\Delta\psi\,\dd s\le\frac{(n-1)r^2}2 .
\]
Hence
\[
0\le\operatorname{avg}\phi-\phi(x_i)\le-\sum_{j\neq i}\log\frac{\max(r,d_j)}{d_j}+\frac{(n-1)r^2}{2}.
\]

\emph{Chart separation.} Each logarithm is $\ge0$, and it is $>\log2$ if $d_j<r/2$. Take
$r^2=2\log2/(n-1)$: no $d_j<r/2$ is possible, so
\[
d_j\ge r/2=\sqrt{\log2/(2(n-1))} .
\]

\emph{Chordal separation.} If $|x_j|\le3$, then $\lambda(x_j)\ge1/5$ and $\lambda(x_i)\ge1$, so
\[
\|y_i-y_j\|\ge5^{-1/2}d_j\ge\sqrt{\log2/(10n)} .
\]
If $|x_j|>3$, then the third coordinates satisfy $(y_j)_3>0.8$ and $(y_i)_3\le0$, so
$\|y_i-y_j\|>0.8$.

\emph{Other charts and blow-up.} For any chart, $|x_a-x_b|=\|y_a-y_b\|/\sqrt{\lambda_a\lambda_b}\ge\|y_a-y_b\|/2$.
Then $|z_a-z_b|\ge\kappa c_{\rm sep}/2=1$.
\end{proof}

\begin{remark}[Theorem 5.2 of the Lectures and a misprint in its statement]\label{rem:E2}
The argument above is the isotropic averaging of Theorem~5.2 of the Lectures, done in
the chart with explicit constants. The statement (5.3.4) there has
$\|(-\Delta V)_+\|_{L^\infty}$. Its hypothesis, its proof (5.3.5) and
its choice of $\eta$ all use $(\Delta V)_+$.
\begin{itemize}[leftmargin=*]
\item This reading confirms a misprint in (5.3.4) (see the errata record of the repository).
\item For $V=\log(1+|x|^2)$, $\Delta V>0$, so the printed statement would give no bound.
\item The proof here does not depend on E2.
\end{itemize}
\end{remark}

%======================================================================
\section{Upper bound: latitude-ring configurations}\label{sec:UB}

\subsection{Row configurations and the exact energy identity}

\begin{definition}[Row configurations]\label{def:rows}
A \emph{row configuration} consists of:
\begin{itemize}[leftmargin=*]
\item heights $t_1<\dots<t_K$;
\item counts $N_k\ge1$ with $\sum_kN_k=n$;
\item phases $\varphi_k$.
\end{itemize}
Its points are $\mathcal M(\varphi_k+2\pi a/N_k,\,t_k)$ for $0\le a<N_k$ and $1\le k\le K$.
\end{definition}

\begin{lemma}[Row sums]\label{lem:rowsums}
Let $\Theta(s):=-\log(1-e^{-s})$ for $s>0$. For a row configuration:
\begin{enumerate}[label=(\roman*)]
\item The ordered pairs inside row $k$ contribute $-N_k\log N_k$ to
  $\sum_{i\ne j}-\log|2\sinh(w_{ij}/2)|$.
\item For $k\ne l$, put $D=|t_k-t_l|$, $\gamma=\gcd(N_k,N_l)$ and $L=\operatorname{lcm}(N_k,N_l)$.
  The pairs $(a,b)$ with $a$ in row $k$ and $b$ in row $l$ contribute
  \[
  -\tfrac12N_kN_lD+X_{kl},\qquad
  X_{kl}:=-\gamma\log\bigl|1-e^{-LD}e^{iL(\varphi_l-\varphi_k)}\bigr| .
  \]
  Moreover $X_{kl}=X_{lk}$.
\item $|X_{kl}|\le\min(N_k,N_l)\,\Theta\bigl(\max(N_k,N_l)D\bigr)$. The function $\Theta$ is
  decreasing, satisfies $\Theta(s+a)\le e^{-a}\Theta(s)$ for $a\ge0$, and satisfies
  $\Theta(s)\le2e^{-s}$ for $s\ge\log2$.
\end{enumerate}
\end{lemma}
\begin{proof}
\emph{(i)} $\prod_{j=1}^{N-1}2\sin(\pi j/N)=N$ (evaluate $\prod_{j=1}^{N-1}(z-\omega^j)$ at $z=1$).

\emph{(ii)} For $\mathrm{Re}\,w=D>0$ we have $|2\sinh(w/2)|=e^{D/2}|1-e^{-w}|$. For each
fixed point $b$ of row $l$, the product over the $N_k$-th roots of unity gives
$\prod_a(1-\zeta\omega^a)=1-\zeta^{N_k}$. As $b$ runs over row $l$, the factor
$e^{2\pi iN_kb/N_l}$ runs over the $(N_l/\gamma)$-th roots of unity, each $\gamma$ times. Apply
the product formula again.

\emph{(iii)} The phases lie on the unit circle, so $|X_{kl}|\le\gamma\Theta(LD)$, with $\gamma\le\min$ and $L\ge\max$.
Put $u=e^{-s}$. Then
$(\log\Theta)'=-\frac{u}{1-u}\big/\bigl(-\log(1-u)\bigr)\le-1$, because $-\log(1-u)\le u/(1-u)$.

\emph{Checks.} (i)--(ii) agree with brute force to $10^{-8}$ absolute (about $10^{-15}$ relative)
for $n\le6001$, for row configurations of the type of Construction~\ref{con:rings}, and for
Construction~\ref{con:rings} itself for $n\le20000$ ($|E_{\rm rows}-E_{\rm brute}|\le1.5\cdot10^{-8}$).
\end{proof}

\begin{definition}[Slab-barycentric row configurations]\label{def:slab}
Let
\[
\nu(\dd t)=\varrho(t)\dd t,\qquad \varrho(t)=\tfrac n2\operatorname{sech}^2t,\qquad
\Phi(t)=\nu((-\infty,t])=\tfrac n2(1+\tanh t).
\]
Then $\nu(\R)=n$ and $\Phi$ has inverse $\tau(a)=\tfrac12\log\frac a{n-a}$. Given counts $N_k$ with
partial sums $A_k=N_1+\dots+N_k$ ($A_0=0$, $A_K=n$), set:
\begin{itemize}[leftmargin=*]
\item slabs $I_k=[\tau(A_{k-1}),\tau(A_k)]$, with $\tau(0)=-\infty$ and $\tau(n)=+\infty$;
\item barycentres $t_k=N_k^{-1}\int_{I_k}t\dd\nu$;
\item local aspect $x_k:=N_k^2/\varrho(t_k)$;
\item cell term
  \[
  \mathrm{Cell}_k:=\int_{I_k}\Bigl(N_k\mathbf 1_{t\ge t_k}-\nu([\tau(A_{k-1}),t])\Bigr)^2\dd t ;
  \]
\item $B_k:=-\tfrac{N_k}2\log x_k+\mathrm{Cell}_k+\sum_{l\ne k}X_{kl}$.
\end{itemize}
Phases are arbitrary.
\end{definition}

\begin{lemma}[1D continuum facts]\label{lem:1d}
Let $U_\nu(t):=\int-\tfrac12|t-s|\dd\nu(s)$. Then:
\begin{enumerate}[label=(\roman*)]
\item $U_\nu(t)=-\tfrac n2\log(2\cosh t)$.
\item $\int\log\cosh\dd\nu=n(1-\log2)$.
\item With
  \[
  \mathcal J(\nu'):=-\tfrac12\iint|t-s|\dd\nu'\dd\nu'+(n-1)\int\log\cosh\dd\nu',
  \]
  we have $\mathcal J(\nu)=(\tfrac12-\log2)n^2-n(1-\log2)$.
\item For every signed measure $\sigma$ with $\sigma(\R)=0$ and $\int|t|\dd|\sigma|<\infty$,
  \[
  -\tfrac12\iint|t-s|\dd\sigma\dd\sigma=\int_\R F_\sigma(u)^2\dd u,\qquad F_\sigma(u)=\sigma((-\infty,u]).
  \]
\end{enumerate}
\end{lemma}
\begin{proof}
\emph{(i)} $U_\nu''=-\varrho=-(\tfrac n2\log\cosh)''$. Both $U_\nu$ and
$-\tfrac n2\log(2\cosh t)$ are even and behave like $-\tfrac n2|t|+o(1)$ at $\pm\infty$, so they
coincide.

\emph{(ii)} Substitute $u=\tanh t$:
$\int\log\cosh\dd\nu=\tfrac n2\int_{-1}^1-\tfrac12\log(1-u^2)\dd u=n(1-\log2)$.

\emph{(iii)}
$\mathcal J(\nu)=\int U_\nu\dd\nu+(n-1)n(1-\log2)=-\tfrac{n^2}{2}(\log2+1-\log2)+(n^2-n)(1-\log2)$.

\emph{(iv)} Write $|t-s|=\int_\R(\mathbf 1_{s\le u}-\mathbf 1_{t\le u})^2\dd u$. Fubini applies by the
moment assumption, and $\sigma(\R)=0$ makes the inner double integral equal to $-2F_\sigma(u)^2$.
\end{proof}

\begin{proposition}[Exact ring identity]\label{prop:ring-identity}
For every slab-barycentric row configuration (any phases),
\[
E=\Bigl(\tfrac12-\log2\Bigr)n^2-\tfrac n2\log n+\tfrac n2\log2+\sum_{k=1}^KB_k .
\]
\end{proposition}
\begin{proof}
\emph{Row decomposition.} By Lemmas~\ref{lem:chordal} and~\ref{lem:rowsums},
$E=E_{1D}+\sum_k\sum_{l\ne k}X_{kl}$ with
\[
E_{1D}=-\sum_kN_k\log N_k+\mathcal J(\nu_d),\qquad \nu_d=\sum_kN_k\delta_{t_k}.
\]
The diagonal does not contribute, because $|t-t|=0$.

\emph{Expansion around $\nu$.} Let $\sigma=\nu_d-\nu$, so $\sigma(\R)=0$. Since
$2U_\nu+(n-1)\log\cosh=-n\log2-\log\cosh$, expanding the quadratic gives
\[
\mathcal J(\nu_d)=\mathcal J(\nu)-\int\log\cosh\dd\sigma-\tfrac12\iint|t-s|\dd\sigma\dd\sigma .
\]

\emph{The cell term.} $F_\sigma$ vanishes at every slab endpoint, because $t_k$ lies in the
interior of $I_k$. On $I_k$ it equals $N_k\mathbf 1_{t\ge t_k}-\nu([\tau(A_{k-1}),t])$. So by
Lemma~\ref{lem:1d}(iv) the last term is $\sum_k\mathrm{Cell}_k$.

\emph{The Jensen term.} Set $\mathrm{Jen}_k:=\int_{I_k}\log\cosh\dd\nu-N_k\log\cosh t_k$. Then
$-\int\log\cosh\dd\sigma=\sum_k\mathrm{Jen}_k$.

\emph{Splitting $N_k\log N_k$.} Write
$-N_k\log N_k=-\tfrac{N_k}2\log x_k-\tfrac{N_k}{2}\log\varrho(t_k)$ and
$-\log\varrho(t_k)=-\log\tfrac n2+2\log\cosh t_k$. Using Lemma~\ref{lem:1d}(ii),
\[
-\sum_k\tfrac{N_k}2\log\varrho(t_k)=-\tfrac n2\log\tfrac n2+n(1-\log2)-\sum_k\mathrm{Jen}_k .
\]
The Jensen terms cancel exactly. Lemma~\ref{lem:1d}(iii) then gives the claim.

\emph{Check.} Verified against brute force to $\le1.4\cdot10^{-8}$ absolute for
$n\in\{40,101,500,1201,3000,6001\}$ (the identity holds for every slab-barycentric row
configuration, so any such configuration may be used for the check). For Construction~\ref{con:rings}
as written, $|E_{\rm id}-E_{\rm brute}|\le3\cdot10^{-8}$ (relative $\le6\cdot10^{-16}$) for
$n\le20000$.
\end{proof}

\begin{remark}
Proposition~\ref{prop:ring-identity} reduces the upper bound to
$\sum_kB_k\le n\,e(\sqrt3\pi)+o(n)$, because $e(\sqrt3\pi)+\tfrac12\log2=C_{\rm ring}$. There is
no continuum Riemann sum left to control: rounding enters only through the $x_k$.
\end{remark}

\subsection{The construction}

\begin{construction}[Ring configuration $Y_n$]\label{con:rings}
\emph{Parameters.} Let $n\ge n_0$, $T:=\lceil n^{1/4}\rceil$ and $N_1^{\rm cap}:=\max(40,4T)$. Put
\[
N_{\rm eq}=\sqrt{\sqrt3\pi n/2},\qquad N(t)=N_{\rm eq}\operatorname{sech}t,\qquad
N^*(a):=N(\tau(a))=\sqrt{2\sqrt3\pi\,a(n-a)/n}.
\]

\emph{Recursion.} Define counts recursively. At step $k$, let $a=A_{k-1}$, $r=n-a$ and
$\tilde n=\max(3,\lfloor N^*(a)+\tfrac12\rfloor)$. Call step $k$ \emph{bulk} if
$N^*(a)\ge N_1^{\rm cap}$.
\begin{enumerate}[label=(\Alph*)]
\item[(F)] If $r\le2\tilde n+3$: set $N_k=r$ and $K=k$, and stop.
\item[(C)] Else, if $k\ge2$, both steps $k$ and $k-1$ are bulk, and $|N_{k-1}-N^*(a)|\le T$: set
  $N_k=N_{k-1}$ and $j_k=j_{k-1}+1$ (continue the band).
\item[(S)] Otherwise set $N_k=\tilde n$ and $j_k=0$ (start a new band).
\end{enumerate}
At step $k=1$, rule (S) gives $N_1=3$.

\emph{Geometry and phases.} Rows, slabs and barycentres are as in Definition~\ref{def:slab}, with
phase $\varphi_k=j_k\pi/N_k$. A \emph{band} is a maximal run of consecutive bulk rows: one row made
by (S), followed by rows made by (C).
\end{construction}

The construction was also implemented numerically. A variant with $T=\lceil n^{1/8}\rceil$ and $N_1=4T$
agrees with Construction~\ref{con:rings} for $n\ge10^4$, where $4T\ge40$. On
Construction~\ref{con:rings} as written, the last count $N_K$ reaches $49$ (at $n=12345$),
within the bound $N_K\le51$ of (G1).

\begin{lemma}[Geometry of the construction]\label{lem:geometry}
There are absolute constants $c_0,C_0>0$ and $n_0$ such that for $n\ge n_0$:
\begin{enumerate}[label=(G\arabic*)]
\item \emph{Termination and counts.} The recursion terminates, $\sum_kN_k=n$, every $N_k\ge3$,
  and $N_K\le51$. Rule (F) is used only at the last step, and that step is not bulk.
\item \emph{Rows near the target count.} For $2\le k\le K-1$: $|N_k-N^*(A_{k-1})|\le T$ if
  step $k$ is bulk, and $\le\tfrac12$ otherwise. In both cases
  $N_k/N^*(A_{k-1})\in[\tfrac34,\tfrac54]$ and $N_k\le\tfrac12(n-A_{k-1})$.
\item \emph{Slab lengths.} For $2\le k\le K-1$:
  \[
  |I_k|=\tfrac12\log\frac{A_k(n-A_{k-1})}{A_{k-1}(n-A_k)}\in\Bigl[\frac{c_0}{N_k},\frac{C_0}{N_k}\Bigr].
  \]
  In the bulk, $|I_k|\le C_0/N_1^{\rm cap}$.
\item \emph{Barycentre depth.} For every $k$, $t_k$ lies in the interior of $I_k$, at distance
  $\ge c_0/N_k$ from $\partial I_k$.
\item \emph{Comparison with $N(t)$.} For every $k$: $N_k/C_0\le N(t)\le C_0N_k$ for all
  $t\in I_k$ with $|t-t_k|\le1$. Hence $x_k\in[c_0,C_0]$. For bulk rows,
  $x_k\in[\tfrac12\sqrt3\pi,2\sqrt3\pi]$ and $|x_k-\sqrt3\pi|\le C_0(T+1)/N_k$.
\item \emph{Counting.}
  \begin{itemize}[leftmargin=*]
  \item $K\le C_0\sqrt n$;
  \item the number of bands is at most $2N_{\rm eq}/(T-\tfrac12)+1$;
  \item the points in non-bulk rows number at most $C_0(N_1^{\rm cap})^2$.
  \end{itemize}
\end{enumerate}
\end{lemma}
\begin{proof}
\emph{Mass-coordinate formulas.} $\varrho(\tau(a))=\tfrac2na(n-a)$ and $N^*(a)^2=\sqrt3\pi\varrho(\tau(a))$.
Since $\max(a,n-a)/n\in[\tfrac12,1]$,
\[
\frac{N^*(a)^2}{2\sqrt3\pi}\le\min(a,n-a)\le\frac{N^*(a)^2}{\sqrt3\pi} .
\]
The slab length formula is $\int\dd a/\varrho(\tau(a))$, because $\dd t=\dd a/\varrho$.

\emph{(G1).} Each $N_k\ge3$, so the recursion terminates, and $\sum N_k=n$ by construction.
Suppose (F) applies at mass $a$. Then $N^*(a)^2\le2\sqrt3\pi r$, so
$r\le2\tilde n+3\le2\sqrt{2\sqrt3\pi r}+4$, which forces $r\le51$ and $N^*(a)<24<N_1^{\rm cap}$.

At a non-final step, $r>2\tilde n+3$. We claim $N_{k}\le r/2$:
\begin{itemize}[leftmargin=*]
\item in non-bulk steps $N_k=\tilde n$;
\item in bulk steps $N_k\le\tfrac54N^*(a)$ while $r\ge N^*(a)^2/(2\sqrt3\pi)\ge3.6N^*(a)$
  (because $N^*\ge40$).
\end{itemize}
Hence $r-N_k\ge3$ at every non-final step, and $N_K\ge3$.

\emph{(G2).} Rules (C) and (S) give the two bounds directly. Ratios: in the bulk,
$T\le N_1^{\rm cap}/4\le N^*/4$. For $2\le k\le K-1$ we have $a\ge3$ and $n-a\ge9$, so
$N^*(a)\ge4$ and $\tfrac12\le N^*/8$.

\emph{(G3).} Apply $\tfrac{y}{1+y}\le\log(1+y)\le y$ to the two factors
$A_k/A_{k-1}=1+N_k/A_{k-1}$ and $(n-A_{k-1})/(n-A_k)=1+N_k/(n-A_k)$. Then use (G2),
$\min(A_{k-1},n-A_{k-1})\asymp N^*(A_{k-1})^2\asymp N_k^2$, and $n-A_k\ge\tfrac12(n-A_{k-1})$.

\emph{(G4) for interior rows.} For $2\le k\le K-1$, $\Phi$ ranges over $[A_{k-1},A_k]$ on $I_k$, and
$A_k\le A_{k-1}+N_k\le C\,A_{k-1}$ (same for the masses on the other side). Hence
$\varrho=\tfrac2n\Phi(n-\Phi)$ varies on $I_k$ by at most an absolute factor $\Lambda_0$. A
probability density on an interval whose sup/inf ratio is $\le\Lambda_0$ has its barycentre at
distance $\ge|I_k|/(2\Lambda_0)$ from both ends. Combine with (G3).

\emph{(G4) for polar rows.} For $k=1$, $N_1=3$ and, in the variable $a=\Phi(t)\le3$,
\[
\tau(3)-t_1=\tfrac13\int_0^3\tfrac12\log\frac{3(n-a)}{a(n-3)}\dd a=\tfrac12+O(1/n).
\]
The case $k=K$ is symmetric with $3\le N_K\le51$.

\emph{(G5).} Use $N(t)^2=\sqrt3\pi\varrho(t)$, $|(\log N)'|=|\tanh|\le1$, and (G2)--(G4).
\begin{itemize}[leftmargin=*]
\item For interior rows, $N(t_k)/N^*(A_{k-1})=e^{O(|I_k|)}$ and
  $x_k=\sqrt3\pi(N_k/N(t_k))^2$.
\item In the bulk, $|I_k|\le\frac12(N_k/A_{k-1}+N_k/(n-A_k))\le 7/N^*(A_{k-1})+O(1/n)\le 7/N_1^{\rm cap}\to0$.
  With (G2) this gives $N_k/N(t_k)\in[\tfrac34-o(1),\tfrac54+o(1)]$, hence
  $x_k\in[\tfrac12\sqrt3\pi,2\sqrt3\pi]$ for $n\ge n_0$. More precisely,
  $N_k/N(t_k)=1+O((T+1)/N_k)$.
\item For the polar rows, $\varrho(t_k)\asymp N_k$ by the computation above, so
  $x_k\asymp N_k\in[3,51]$.
\end{itemize}

\emph{(G6), rows.} Interior slabs satisfy $\int_{I_k}N(t)\dd t\ge|I_k|N_k/C_0\ge c_0/C_0$, and
$\int_\R N=\pi N_{\rm eq}$. So $K\le2+C\sqrt n$.

\emph{(G6), bands.} Consider a band that starts at row $k_0$ and is ended by an (S)-row $k_1+1$
that is still bulk. Then $|N^*(A_{k_1})-N^*(A_{k_0-1})|>T-\tfrac12$. These mass intervals are
disjoint, and $N^*$ is unimodal on $[0,n]$ with total variation $2N_{\rm eq}$. The bulk steps
form one run, because $\{N^*\ge N_1^{\rm cap}\}$ is an interval. So at most one band ends by
leaving the bulk.

\emph{(G6), non-bulk points.} Non-bulk rows start at masses $a$ with
$\min(a,n-a)\le C(N_1^{\rm cap})^2$, and each holds at most $CN_1^{\rm cap}$ points.
\end{proof}

\subsection{Estimates}

\begin{lemma}[Cell term]\label{lem:cell}
For interior rows ($2\le k\le K-1$), let $\Lambda_k$ be the ratio $\sup/\inf$ of $\varrho$ on $I_k$.
Then
\[
\mathrm{Cell}_k\le\frac{N_kx_k}{12}\,\Lambda_k\bigl(1+3(\Lambda_k-1)^2\bigr)
\quad\text{and}\quad
\mathrm{Cell}_k\le N_k^2|I_k|\le C_0N_k .
\]
For the two polar rows, $\mathrm{Cell}_k\le C$.
\end{lemma}
\begin{proof}
\emph{Change of variable.} Let $G(t)=\nu([\tau(A_{k-1}),t])$ and substitute $\gamma=G(t)$,
$\dd t=\dd\gamma/\varrho$. With $\alpha=G(t_k)$,
\[
\mathrm{Cell}_k\le\frac{\alpha^3+(N-\alpha)^3}{3\varrho_{\min}}
=\frac{N^3/4+3N(\alpha-N/2)^2}{3\varrho_{\min}} .
\]

\emph{Barycentre against median.} Both $t_k-\tau$ and $G^{-1}(N/2)-\tau$ lie in
$[N/(2\varrho_{\max}),N/(2\varrho_{\min})]$. Hence $|\alpha-N/2|\le\tfrac N2(\Lambda_k-1)$.
Finally $1/\varrho_{\min}\le\Lambda_k/\varrho(t_k)$.

\emph{Crude bound.} $|F|\le N_k$ on $I_k$, and use (G3).

\emph{Polar rows.} $F$ is bounded by $N_k\le51$ and decays exponentially beyond $t_k$, because
$\nu((t,\infty))\le\nu((t_k,\infty))e^{-(t-t_k)}$ for $t\ge t_k\ge0$; symmetrically at the south
pole.
\end{proof}

\begin{lemma}[Row-sum bound]\label{lem:rowsum}
There is an absolute $C_X$ such that for $n\ge n_0$, every $k$ and every $u_0\ge0$,
\[
\sum_{l\ne k,\ |t_l-t_k|\ge u_0}\min(N_k,N_l)\,\Theta\bigl(\max(N_k,N_l)|t_l-t_k|\bigr)\le C_XN_ke^{-N_ku_0/2} .
\]
\end{lemma}
\begin{proof}
\emph{Reduction.} Write $M_l=\max(N_k,N_l)$ and use $\min(N_k,N_l)\le N_k$ together with
$\Theta(s+a)\le e^{-a}\Theta(s)$. It suffices to prove
\[
\Sigma:=\sum_{l\neq k}\Theta\bigl(\tfrac12M_l|t_l-t_k|\bigr)\le C_X .
\]

\emph{Separation from row $k$.} By (G4), $t_l\notin I_k^\circ$ and $t_k\notin I_l^\circ$, so
$|t_l-t_k|\ge c_0\max(1/N_k,1/N_l)$.

\emph{Dyadic classes and counting.} Split $l$ into classes $L_j=\{2^j\le N_l<2^{j+1}\}$, $j\ge1$.
By (G3)--(G4), the interior slabs of rows in $L_j$ are disjoint, have length in
$[c_02^{-j-1},C_02^{-j}]$, and contain their barycentres. Hence
\[
\#\{l\in L_j:|t_l-t_k|\le u\}\le\tfrac{4}{c_0}(2^ju+C_0)+2 .
\]

\emph{Distance from the count mismatch.} By (G5) and $|(\log N)'|\le1$,
$|t_l-t_k|\ge\delta_j:=\bigl(|j\log2-\log N_k|-3\log C_0\bigr)_+$ for $l\in L_j$.

\emph{Summation.} By summation by parts in $u$, with $v=\tfrac12M_lu$ and
$M_l\ge\max(N_k,2^j)$, the class $L_j$ contributes at most $\Psi(v_j)$, where:
\begin{itemize}[leftmargin=*]
\item $\Psi(v_*):=C\int_{v_*}^\infty(-\Theta'(v))(1+v)\dd v+C\Theta(v_*)(1+v_*)\le C(1+v_*)e^{-v_*}$;
\item $v_j\ge\max(c_0/4,\tfrac12\max(N_k,2^j)\delta_j)$.
\end{itemize}
For the $O(1)$ classes with $|j-\log_2N_k|\le J_0$ use $\Psi(c_0/4)$. For the other classes,
$\tfrac12\max(N_k,2^j)\delta_j\ge c\,\max(N_k,2^j)\,|j-\log_2N_k|$, and
$\sum_j(1+v_j)e^{-v_j}$ is bounded by an absolute constant (recall $N_k\ge3$). The two polar rows
add at most $2\Theta(c_0/2)$.
\end{proof}

\begin{lemma}[In-band spacing]\label{lem:inband}
Let rows $k$ and $k+m$ ($m\ge1$) lie in the same band, with count $N_b$. Then
\[
\bigl|N_b(t_{k+m}-t_k)-m\,x_k\bigr|\le C\,\frac{m^2}{N_b},\qquad
X_{k,k+m}=-N_b\log\bigl|1-(-1)^me^{-N_b(t_{k+m}-t_k)}\bigr| .
\]
\end{lemma}
\begin{proof}
\emph{The exact formula for $X$.} Here $\gamma=L=N_b$ and $L(\varphi_{k+m}-\varphi_k)=m\pi$;
apply Lemma~\ref{lem:rowsums}(ii).

\emph{Adjacent rows.} As in Lemma~\ref{lem:cell}, $t_j-\tau(A_{j-1})$ and $\tau(A_j)-t_j$ are
$\frac{N_b}{2\varrho(t_j)}(1+O(|I_j|))$. So
\[
N_b(t_{j+1}-t_j)=\tfrac12(x_j+x_{j+1})+O(1/N_b).
\]

\emph{Drift of $x_j$.} $x_j/x_k=\varrho(t_k)/\varrho(t_j)=e^{O(|t_j-t_k|)}$ with
$|t_j-t_k|\le C|j-k|/N_b$, by (G3)--(G4). Sum over $j=k,\dots,k+m-1$.
\end{proof}

\begin{proposition}[Per-row bounds]\label{prop:perrow}
There is an absolute $C$ such that for $n\ge n_0$:
\begin{enumerate}[label=(\roman*)]
\item for every row, $B_k\le CN_k$;
\item for a bulk row $k$ in a band of count $N_b$, let $d_k\ge1$ be the index distance from $k$ to
  the nearest row not in the band (so rows $k\pm m$ with $1\le m<d_k$ are in the band). Then
  \[
  B_k\le N_k\,e(x_k)+C+CN_ke^{-d_k/C}.
  \]
\end{enumerate}
\end{proposition}
\begin{proof}
\emph{(i)} We have $-\tfrac{N_k}{2}\log x_k\le\tfrac{N_k}2\log(1/c_0)$ by (G5), the cell bound of
Lemma~\ref{lem:cell}, and $\sum_l|X_{kl}|\le C_XN_k$ by Lemmas~\ref{lem:rowsums}(iii)
and~\ref{lem:rowsum} with $u_0=0$.

\emph{(ii), cell term.} In the bulk, $\Lambda_k\le e^{2|I_k|}\le1+C/N_k$. Lemma~\ref{lem:cell} and
(G5) give $\mathrm{Cell}_k\le N_kx_k/12+C$.

\emph{(ii), rows inside the band.} Split $\sum_{l\ne k}X_{kl}$ into $|l-k|<d_k$ (rows in the
band, on both sides) and the rest. For the first part, Lemma~\ref{lem:inband} and the bound
$|\frac{\dd}{\dd s}\log|1\mp e^{-s}||\le2e^{-s}$ for $s\ge\log2$ give
\[
\sum_{1\le|l-k|<d_k}X_{kl}\le-2N_b\sum_{m=1}^{d_k-1}\log\bigl|1-(-1)^me^{-mx_k}\bigr|+C\sum_mm^2e^{-mx_k/2}.
\]
(For $m$ beyond $cN_b$ both sides are $O(N_be^{-cm})$.) Completing the sum to $m\ge1$ costs at
most $CN_be^{-x_k(d_k-1)}$.

\emph{(ii), rows outside the band.} These rows have $|t_l-t_k|\ge c_0(d_k-1)/N_b$, because
consecutive rows of the band are $\ge c_0/N_b$ apart by (G4). Lemma~\ref{lem:rowsum} bounds them by
$C_XN_ke^{-c_0(d_k-1)/2}$.

\emph{(ii), conclusion.} Collect the terms:
$-\tfrac{N_k}2\log x_k+\tfrac{N_kx_k}{12}-2N_k\sum_m\log|\cdots|=N_ke(x_k)$.
\end{proof}

\begin{theorem}[Upper bound]\label{thm:UB}
For $n\ge n_0$, the configuration $Y_n$ of Construction~\ref{con:rings} satisfies
\[
E_{\log}(n)\le E(Y_n)\le\Bigl(\tfrac12-\log2\Bigr)n^2-\tfrac n2\log n+C_{\rm ring}\,n+C\,n^{3/4} .
\]
In particular $\limsup_n\min\mathfrak K/n\le C_{\rm ring}$, and $\min\mathfrak K\le C_Un$ for
all $n\ge2$, with an absolute $C_U$.
\end{theorem}
\begin{proof}
\emph{Splitting the rows.} By Proposition~\ref{prop:ring-identity} it suffices to bound
$\sum_kB_k-n\,e(\sqrt3\pi)$. Split the rows into bulk and non-bulk.

\emph{Non-bulk rows.} Proposition~\ref{prop:perrow}(i) and (G6) give
\[
\sum_{\text{non-bulk}}\bigl(B_k-N_ke(\sqrt3\pi)\bigr)\le C(N_1^{\rm cap})^2\le C\sqrt n .
\]

\emph{Bulk rows.} Proposition~\ref{prop:perrow}(ii) applies. On
$[\tfrac12\sqrt3\pi,2\sqrt3\pi]$, $e$ is Lipschitz, and (G5) gives
$N_k|e(x_k)-e(\sqrt3\pi)|\le C(T+1)$. Summing:
\begin{itemize}[leftmargin=*]
\item the Lipschitz term contributes $C(T+1)K\le Cn^{3/4}$;
\item the constant $C$ per row contributes $CK\le C\sqrt n$;
\item the terms $CN_ke^{-d_k/C}$ contribute $\le C'N_b$ per band, hence by (G6) at most
  \[
  C(2N_{\rm eq}/(T-\tfrac12)+1)(N_{\rm eq}+T)\le Cn^{3/4}.
  \]
\end{itemize}

\emph{Constants.} Use $\sum_kN_k=n$ and $e(\sqrt3\pi)+\tfrac12\log2=C_{\rm ring}$. For small $n$
the bound $\min\mathfrak K\le C_Un$ is trivial.

\emph{Numerics.} $c_n-C_{\rm BHS}$ equals $+0.00103$, $+0.00050$ and $+0.00034$ at
$n=10^6$, $10^7$ and $4\cdot10^7+1$ ($T=\lceil n^{1/4}\rceil$).
The value is positive and decreasing, with odd and even $n$ alike.
\end{proof}

\begin{remark}[Stationarity: not needed]\label{rem:stationarity}
$e'(x)=-\frac1{2x}+\frac{E_2(\tau_x)}{12}$ with $\tau_x=\tfrac12+\tfrac{ix}{2\pi}$. Two facts
combine:
\begin{itemize}[leftmargin=*]
\item $E_2$ is $1$-periodic and satisfies $E_2(-1/\tau)=\tau^2E_2(\tau)-6i\tau/\pi$;
\item $\rho-1=-1/\rho$.
\end{itemize}
Together they give $E_2(\rho)=2\sqrt3/\pi$, hence $e'(\sqrt3\pi)=0$. This is the criticality of
the triangular lattice in the centred-rectangular family (numerically $e'(\sqrt3\pi)=-4\cdot10^{-18}$;
the $E_2$ check is $10^{-42}$).

Using it in Proposition~\ref{prop:perrow} improves the in-band error to $O(T^2)$ per band, so the
total becomes $O(n/T+T^2+\sqrt n)$ for $T=o(\sqrt n)$. The proof of Theorem~\ref{thm:UB} does not
use it, which avoids using the quasi-modularity of $E_2$.
\end{remark}

\begin{lemma}[Insertion bound; not needed]\label{lem:insertion}
For $n\ge2$ and $k\ge1$,
\[
E_{\log}(n+k)\le E_{\log}(n)+(1-2\log2)nk+(\tfrac12-\log2)k(k-1).
\]
\end{lemma}
\begin{proof}
Let $\sigma$ be the uniform probability on $S^2$. Then $\int-\log\|y-y'\|\dd\sigma(y')=\tfrac12-\log2$
for every $y$; this follows from Lemma~\ref{lem:U} via $T$. Average the energy of a minimizer for $n$, augmented by $k$ independent
$\sigma$-distributed points. There are $2nk$ ordered old--new pairs and $k(k-1)$ ordered new--new
pairs.
\end{proof}

%======================================================================
\section{Lower bound}\label{sec:LB}

\subsection{Cited input: screening}

\begin{definition}[Neumann puncture energy on a square]\label{def:neumann}
Let $Q=a_0+[0,t]^2$ be a closed square, $\bar m>0$ with $\bar n:=\bar mt^2\in\N$, and
$Y=(y_1,\dots,y_{\bar n})$ pairwise distinct points in $Q^\circ$.
\begin{itemize}[leftmargin=*]
\item \emph{Density one} ($\bar m=1$, so $t=\sqrt{\bar n}$). Reflect $Y$ evenly across the sides
  of $Q$. This gives $4\bar n$ pairwise distinct classes $\tilde Y$ modulo
  $\Gamma:=2t\,\Z[i]=\Gamma_{4\bar n}$. Let
  $\tilde h=\mathtt{torusPotential}(\Gamma,\tilde Y)$ and define $u_Y:=\tilde h|_Q$ and
  $W^{\rm Neu}_Q(Y;1):=W_{Q^\circ}(u_Y)$.
\item \emph{General $\bar m$.} Put $\sigma=\sqrt{\bar m}$, $u_Y(x):=u_{\sigma Y}(\sigma x)$ and
  $W^{\rm Neu}_Q(Y;\bar m):=W_{Q^\circ}(u_Y)$.
\end{itemize}
\end{definition}

\begin{lemma}[Reflection and Theorem 1.2 of the paper]\label{lem:reflection}
In Definition~\ref{def:neumann}:
\begin{enumerate}[label=(\roman*)]
\item $\tilde h$ is invariant under the reflections across the lines bounding $Q$ (and their
  translates by $\Gamma$).
\item $u_Y\in C^\infty(Q\setminus Y)$ solves $-\Delta u_Y=2\pi(\sum\delta_{y_i}-\bar m)$ in $Q^\circ$
  with $\partial_\nu u_Y=0$ on $\partial Q$ (classical Neumann problem, (7.1.1) of the Lectures).
\item $\mathtt{torusEnergy}(\Gamma,\tilde Y)=4W^{\rm Neu}_Q(Y;1)$.
\item $W^{\rm Neu}_Q(Y;\bar m)=W^{\rm Neu}_{\sigma Q}(\sigma Y;1)-\frac\pi2\bar n\log\bar m\ge\bar n\bigl(\Wtri-\frac\pi2\log\bar m\bigr)$.
\end{enumerate}
\end{lemma}
\begin{proof}
\emph{(i)} Let $\mathcal R$ be one of the reflections. It maps $\Gamma$ to itself, preserves
the set $\tilde Y$ mod $\Gamma$ and preserves Lebesgue measure on a fundamental domain. Hence
$\tilde h\circ\mathcal R$ satisfies every clause of $\mathtt{IsTorusPotential}(\Gamma,\tilde Y)$
(the clauses depend only on the set of charges). By uniqueness of the torus potential,
$\tilde h\circ\mathcal R=\tilde h$.

\emph{(ii)} Smoothness off the charges and $\Delta\tilde h=2\pi$ are clauses of the definition.
Evenness across a side forces the normal derivative to vanish on that side, including at the
corners, where both reflections apply.

\emph{(iii)} The square $[-t,t]^2$ (shifted to contain $Q$ as a quarter) is a fundamental domain.
It is the union of four reflected copies of $Q$ that overlap in a null set. For small $\eta$ the
$4\bar n$ discs lie inside the open copies. Use (i) and the translation invariance of the torus
energy over fundamental domains (\lean{torusDomain\_spec}).

\emph{(iv)} The scaling identity is Lemma~\ref{lem:conformal} with $k=\sigma$, since the scaled
function solves the density-one problem. The inequality is (iii) together with Theorem~1.2 of the paper
(Input~\ref{in:R1}, $N=4\bar n\ge4$).
\end{proof}

\begin{citedinput}[Screening: Prop.~7.18 of the Lectures with Remark~7.19, specialized; filler by Lemmas~\ref{lem:filler} and~\ref{lem:filler-energy}]\label{in:SCR}
\textbf{Sources:} Prop.~7.18 and Remark~7.19 of the Lectures; the proof is in Appendix~A. Cor.~7.11, which one could use for the
filler, is replaced by Lemmas~\ref{lem:filler}--\ref{lem:filler-energy}. This input is also proved by the explicit screening
of the screening note (Proposition~7.2 there). Specialization:
\begin{itemize}[leftmargin=*]
\item $d=2$, $s=0$, $c_d=2\pi$;
\item $U=\R^2$, inner case;
\item $\Omega=\Omega'=Q$ and $\Omega''=Q^{1/2}$, where $Q^{s}:=\{x\in Q:\operatorname{dist}(x,\partial Q)\ge s\}$;
\item $\mu\equiv\bar m$ constant and $\eta=1$;
\item converted to puncture normalization by Remark~\ref{rem:book-F}.
\end{itemize}
\textbf{Statement.} Let $0<m_-\le m_+\le1/150$. There are $C_S,C_{\rm scr}\ge1$ and $\ell_0\ge1$,
depending only on $(m_-,m_+)$, such that the following holds. Here $\ell_0=\ell_0(m_-,m_+)$ is the
constant $C(d,m)$ of Appendix~A, which the proof uses in the form ``$\ell\ge C$'' and ``$\ell>C>2\eta$'';
Prop.~7.18 as printed does not state
this condition. Let:
\begin{itemize}[leftmargin=*]
\item $Q$ be a closed square of side $t$, and $\bar m\in[m_-,m_+]$ with $\bar n=\bar mt^2\in\N$;
\item $\max(4m_+/m_-,\ell_0)\le\ell\le\tilde\ell\le t/8$;
\item $X=(x_1,\dots,x_p)\subset Q$ with $|x_i-x_j|\ge1$;
\item $w$ such that $w^\natural:=w-\sum_if_{1/4}(\cdot-x_i)\in H^1(Q)$ and
  $-\Delta w=2\pi(\sum_{x_i\in Q^{1/2}}\delta_{x_i}-\bar m)$ in $\mathcal D'((Q^{1/2})^\circ)$.
\end{itemize}
Put $S:=\int_{Q^{\tilde\ell/2}\setminus Q^{2\tilde\ell}}|\nabla w^\natural|^2$. If
$S\le\ell^3\tilde\ell/C_{\rm scr}$, then there are $\bar n$ pairwise distinct points
$Y\subset Q^\circ$ with
\[
W^{\rm Neu}_Q(Y;\bar m)\le\tfrac12\int_Q|\nabla w^\natural|^2+\pi p\log\tfrac14
-2\pi\bar m\sum_{i=1}^p\int_Qf_{1/4}(x-x_i)\dd x
+C_S\Bigl(\frac{\ell S}{\tilde\ell}+t\tilde\ell+|\bar n-p|\Bigr).
\]
\textbf{Translation from the Lectures.}
\begin{itemize}[leftmargin=*]
\item Separation $\ge1$ makes every radius (7.1.24) equal to $\tfrac14$, so $w^\natural=w_{\hat r}$.
\item $S$ dominates $S(X_n,w)$ of (7.2.6) of the Lectures under either reading of $Q_{R-\tilde\ell}$
  (shrink of the side, or of the depth, by $\tilde\ell$). The screenability condition (7.2.8)
  follows.
\item (7.2.14) holds with $\eta=1$. For constants uniform in $\bar m\in[m_-,m_+]$ we take the
  $m:=m_-$ of the Lectures and use $\|\mu\|_{L^\infty}=\bar m\le m_+$; then (7.2.14), $\eta\le\ell m/(4\|\mu\|_{L^\infty})$
  holds with $\eta=1$ as soon as $\ell\ge4m_+/m_-$. In
  Definition~\ref{def:grid}, $m_+/m_-=25$, so this is $\ell\ge100$. The index set $J$ is then empty
  (Remark~7.19). Appendix~A additionally needs $\ell\ge\ell_0(m_-,m_+)$, which is the extra hypothesis
  of the statement above.
\item We take $C_{\rm scr}>C$, where $C$ is the constant in the ``Moreover'' clause of Prop.~7.18. Then $S\le\ell^3\tilde\ell/C_{\rm scr}$ gives
  $\ell^3>CS/\tilde\ell$, also when $S=0$. So $O=\{x\in Q:\operatorname{dist}(x,\partial Q)\ge w\}$ is a
  concentric closed square, with $\tilde\ell/2\le w\le2\tilde\ell$ under either reading of
  $Q_{R-2\tilde\ell}\subset O\subset Q_{R-\tilde\ell}$. Hence $N_\eta$ is a square annulus with rounded
  inner corners.
\item The filler $Z$ is chosen by Lemmas~\ref{lem:filler} and~\ref{lem:filler-energy}: it gives $\bar n-n_O$ distinct points in the
  interior of $N_\eta$ with $F^\flat(Z,\tilde\mu,N_\eta)\le C_F|N_\eta|\le 8C_Ft\tilde\ell$, where $F^\flat$
  is the variational energy of \S\ref{sec:filler}, which satisfies $F\le F^\flat$ (Lemma~7.5 of the Lectures) and is
  provided by (X3') of Remark~\ref{rem:wp13-export} in place of $F$. Lemma~\ref{lem:filler-energy}
  needs only $\tilde\mu\le\bar m+m_-/2\le\tfrac32m_+\le\tfrac1{100}$ on $N_\eta$ and
  $\tilde\mu(N_\eta)=\bar n-n_O$, which are (7.2.15)--(7.2.16). It replaces the appeal of the Lectures to
  (7.1.57) at this step, whose hypotheses (Prop.~7.10/Cor.~7.11) are not
  verified for the thin layer $N_\eta$.
\item $2\pi F(Y,\bar m,Q)=W^{\rm Neu}_Q(Y;\bar m)$ by Lemmas~\ref{lem:truncation}
  and~\ref{lem:reflection}(ii): the radii (7.1.21) give disjoint discs in $Q^\circ$, and the
  classical Neumann solution is unique up to constants.
\end{itemize}
\end{citedinput}

\begin{remark}[Obligations on the proof of Prop.~7.18 for Input~\ref{in:SCR}]\label{rem:scr-interface}
\begin{enumerate}[label=(\alph*)]
\item The filler bound $\min F^\flat(\cdot,\tilde\mu,N_\eta)\le C|N_\eta|$ on
  the thin layer $N_\eta$ is Lemma~\ref{lem:filler-energy}, proved in \S\ref{sec:filler} without
  Proposition~7.10/Corollary~7.11 (whose hypotheses are not checked for $N_\eta$ in the Lectures, which use that step in the proof of Prop.~7.18).
  What remains is an \emph{interface}:
  the proof of Prop.~7.18 must provide the items (X1)--(X3') of Remark~\ref{rem:wp13-export}.
\item Constants uniform for constant densities $\bar m\in[m_-,m_+]$. The constants of the Lectures depend on
  $(d,m,\|\mu\|_\infty)$; monotone dependence must be stated.
\item The output $Y$ lies in $Q^\circ$ and is pairwise distinct. In the Lectures, $Y$ is the kept
  points within $\tfrac14$ of $O\subset Q_{R-\tilde\ell}$, together with the filler in $N_\eta$.
  The filler of Lemma~\ref{lem:filler} lies in the interior of $N_\eta$, so at distance $>1$ from $O$
  and $\ge\tfrac34$ from every kept point; its points are $\tfrac73$-separated, and its Neumann
  energy is bounded by $C_F|N_\eta|$.
\item Prop.~7.18 is proved through (7.1.37), which contains an error (see the errata record of the repository). Use the repaired
  version; only constants change. The same holds
  for (X2)/(X3): they come from Appendix~A with the corrected equation
  $-\Delta h_{2,k}=c_d\bigl(m_k\mathbf 1_{H_k^\eta}-\mu\mathbf 1_{H_k\setminus H_k^\eta}\bigr)$
  (see the errata record), which matters here because $\eta=1>0$.
\item The Neumann potential of the Lectures is identified with $u_Y$ of Definition~\ref{def:neumann}.
\end{enumerate}
\end{remark}

\subsection{Elementary lemmas}

\begin{lemma}[Counting separated points]\label{lem:count}
If $|z_i-z_j|\ge1$ for $i\ne j$, then for every measurable $A$,
$\#\{i:z_i\in A\}\le\frac4\pi|A+B(0,\tfrac12)|$.
\end{lemma}
\begin{proof}
The discs $B(z_i,\tfrac12)$ are disjoint and contained in $A+B(0,\tfrac12)$.
\end{proof}

\begin{lemma}[Density freeze]\label{lem:freeze}
Let $Q$ be a square of side $t$ and $m\in C^1(Q)$ with $\bar m=\frac1{|Q|}\int_Qm$ and
$\delta:=\sup_Q|m-\bar m|$. Then there is $v\in H^1(Q)$ with $\int_Q\nabla v\cdot\nabla\phi=2\pi\int_Q(\bar m-m)\phi$ for all
$\phi\in H^1(Q)$ (weak Neumann problem), and
\[
\int_Q|\nabla v|^2\le4t^4\delta^2 .
\]
\end{lemma}
\begin{proof}
\emph{Existence.} Use the Riesz representation on mean-zero $H^1(Q)$, which requires the Neumann
Poincaré inequality $\|\phi-\bar\phi\|_{L^2(Q)}\le\frac t\pi\|\nabla\phi\|_{L^2(Q)}$. The first
nonzero Neumann eigenvalue of the square is $\pi^2/t^2$ (cosine series); this is a standard fact.

\emph{Estimate.} Test with $\phi=v-\bar v$:
\[
\|\nabla v\|_2^2\le2\pi\|v-\bar v\|_2\|m-\bar m\|_2\le2t\|\nabla v\|_2\cdot t\delta .
\]
\end{proof}

\begin{lemma}[Discrepancy]\label{lem:discrepancy}
Let $H$ be admissible on $(\C,Z)$ with $-\Delta H=2\pi(\sum_i\delta_{z_i}-m)$, $0\le m\le m_+$, and
$|z_i-z_j|\ge1$. Let $H_{1/4}=H-\sum_if_{1/4}(\cdot-z_i)$, and let $Q$ be a square of side
$t\ge2$ with $p$ points of $Z$. Then
\[
|p-m(Q)|\le C(1+m_+)t+\sqrt{2t}\Bigl(\int_{Q^{1/2}\setminus Q^{1}}|\nabla H_{1/4}|^2\Bigr)^{1/2}
\]
with an absolute $C$.
\end{lemma}
\begin{proof}
\emph{Flux identity.} $H_{1/4}$ is locally Lipschitz, and
$-\Delta H_{1/4}=2\pi(\sum_i\sigma_i-m)$, where $\sigma_i$ is the uniform probability on
$\partial B(z_i,\tfrac14)$. For a.e. $s\in[\tfrac12,1]$ the divergence theorem on $Q^s$ gives
\[
\Bigl|\sum_i\sigma_i(Q^s)-m(Q^s)\Bigr|\le\frac1{2\pi}\int_{\partial Q^s}|\nabla H_{1/4}| .
\]

\emph{Averaging in $s$.} Average over $s\in[\frac12,1]$ and apply Cauchy--Schwarz on the layer,
whose area is $\le2t$. Then compare:
\begin{itemize}[leftmargin=*]
\item $|p-\sum_i\sigma_i(Q^s)|$ is at most the number of points within $\tfrac54$ of
  $\partial Q$, which is $\le Ct$ by Lemma~\ref{lem:count};
\item $|m(Q)-m(Q^s)|\le4m_+t$.\qedhere
\end{itemize}
\end{proof}

\subsection{The hemisphere functional}

\begin{definition}[Setup for a minimizer]\label{def:LBsetup}
Let $Y$ minimize $E$ on $(S^2)^n$, rotated so that no point is at $\pm\mathbf n$
(Lemma~\ref{lem:existence}). Let $x_i=T^{-1}(y_i)$, $\tilde x_i=\iota(x_i)$ and use the notation of
Definition~\ref{def:blowup}. For $r\in[1,2]$ put:
\begin{itemize}[leftmargin=*]
\item $d_i(r)=s_n\bigl||x_i|-r\bigr|$ and $\tilde d_i(r)=s_n\bigl||\tilde x_i|-1/r\bigr|$;
\item $\Lambda(r):=\sum_i\bigl[\log_+\frac1{2d_i(r)}+\log_+\frac1{2\tilde d_i(r)}\bigr]$.
\end{itemize}
\end{definition}

\begin{lemma}[Choice of the splitting circle]\label{lem:choose-r}
There is $r\in[1,2]$ with $|x_i|\ne r$ for all $i$ and $\Lambda(r)\le20\sqrt n/\kappa$.
\end{lemma}
\begin{proof}
\emph{First sum.} For fixed $\rho_0$,
\[
\int_1^2\log_+\frac{1}{2s_n|\rho_0-r|}\dd r\le2\int_0^{1/(2s_n)}\log\frac1{2s_nu}\dd u=\frac1{s_n}.
\]

\emph{Second sum.} Substitute $u=1/r$, $\dd r=\dd u/u^2\le4\dd u$; the bound becomes $4/s_n$.

\emph{Conclusion.} Summing, $\int_1^2\Lambda\le5n/s_n=5\sqrt n/\kappa$. Exclude the finite set
$\{|x_i|\}$ and take $r$ with $\Lambda(r)$ at most four times the average.
\end{proof}

\begin{lemma}[Crude bound and a priori bound]\label{lem:apriori}
Let $D=B(0,rs_n)$ and $\tilde D=B(0,s_n/r)$, with $r$ from Lemma~\ref{lem:choose-r}. Let
$\Phi_{\mathcal S},\Phi_{\mathcal N}$ be as in Proposition~\ref{prop:twochart}, and put
$\alpha_i:=\min(\tfrac14,\tfrac12d_i(r))$ for $z_i\in D$. Then:
\begin{enumerate}[label=(\roman*)]
\item
\[
\Phi_{\mathcal S}=\tfrac12\int_D|\nabla H_\alpha|^2+\sum_{z_i\in D}\Bigl[\pi\log\alpha_i-2\pi\int f_{\alpha_i}(\cdot-z_i)m+\pi\log\lambda(x_i)\Bigr],
\]
and
\[
\Phi_{\mathcal S}\ge\tfrac12\int_D|\nabla H_\alpha|^2-C_1n_{\mathcal S}-\pi\Lambda(r).
\]
The same holds for $\Phi_{\mathcal N}$, with $\tilde H$, $\tilde D$ and $\tilde d_i$.
\item With $C_U$ from Theorem~\ref{thm:UB},
\[
\int_D|\nabla H_\alpha|^2+\int_{\tilde D}|\nabla\tilde H_{\tilde\alpha}|^2\le C_An,
\qquad C_A:=2\pi C_U+2\pi|\log\kappa|+4C_1+80\pi/\kappa .
\]
\end{enumerate}
\end{lemma}
\begin{proof}
\emph{(i)} The discs $\bar B(z_i,\alpha_i)$ lie in $D$ and are disjoint, by
Theorem~\ref{thm:separation}. Apply Lemma~\ref{lem:truncation} with $\varrho=m\le1/(\pi\kappa^2)$.
Then use $\lambda\ge2/5$ on $|x|\le2$, $\int f_\alpha\le\pi/32$, and
$\log\alpha_i\ge\log\tfrac14-\log_+\frac1{2d_i}$.

\emph{(ii)} Add the two inequalities of (i) and use Proposition~\ref{prop:twochart}:
\[
\Phi_{\mathcal S}+\Phi_{\mathcal N}=\pi\mathfrak K(Y)+\pi n\log\kappa\le\pi C_Un+\pi n\log\kappa .
\]
\end{proof}

\subsection{Localization and the per-square estimate}

\begin{definition}[Grid, quantized squares]\label{def:grid}
Let
\[
m_-:=\frac{(2/5)^2}{4\pi\kappa^2}\le m\le m_+:=\frac1{\pi\kappa^2}\quad\text{on } D ,
\]
so $m_+/m_-=25$ and $m_+=\log2/(40\pi)<1/150$. Fix $R\ge R_0$, where the absolute constant $R_0$
satisfies $R_0\ge10^5$ (so that $\ell\ge100=4m_+/m_-$, as Input~\ref{in:SCR} requires),
$R_0\ge\ell_0^{5/2}$ (so that $\ell=R^{2/5}\ge\ell_0=\ell_0(m_-,m_+)$, as Input~\ref{in:SCR} requires;
$m_\pm$ are absolute, so $\ell_0$ is an absolute constant) and the other
largeness conditions used below. Fix $\ell=R^{2/5}$ and $\tilde\ell=R^{1/2}$.
\begin{itemize}[leftmargin=*]
\item For an offset $\theta\in[0,R)^2$, let $\mathcal K_\theta$ be the set of squares
  $K=\theta+R(j+[0,1)^2)$, $j\in\Z^2$, with $\bar K\subset D_1:=\{z\in D:\operatorname{dist}(z,\partial D)\ge1\}$.
\item For $K$ with lower-left corner $a_K$, let $Q_K:=a_K+[0,t_K]^2$, with $t_K\in[R-1,R-\tfrac12]$
  chosen so that $m(Q_K)\in\N$. This is possible because $t\mapsto m(a_K+[0,t]^2)$ is continuous
  and increases by $\ge m_-(R-\tfrac34)\ge1$ on that interval.
\end{itemize}
Then the $Q_K$ are pairwise disjoint and $Q_K\subset K$. Set:
\begin{itemize}[leftmargin=*]
\item $\bar n_K:=m(Q_K)$, $\bar m_K:=\bar n_K/t_K^2\in[m_-,m_+]$, and $p_K:=\#(Z\cap Q_K)$;
\item $\mathrm{strip}^+(K):=\{z\in K:\tilde\ell/2\le\operatorname{dist}(z,\partial K)\le2\tilde\ell+1\}$;
\item $\mathrm{lay}(K):=\{z\in K:\operatorname{dist}(z,\partial K)\le2\}$.
\end{itemize}
For $z\in Q_K$, $\operatorname{dist}(z,\partial Q_K)\le\operatorname{dist}(z,\partial K)\le\operatorname{dist}(z,\partial Q_K)+1$.
Hence $Q_K^{\tilde\ell/2}\setminus Q_K^{2\tilde\ell}\subset\mathrm{strip}^+(K)$ (a point with
$\operatorname{dist}(z,\partial Q_K)\in[\tilde\ell/2,2\tilde\ell)$ has $\operatorname{dist}(z,\partial K)\in[\tilde\ell/2,2\tilde\ell+1)$) and
$Q_K^{1/2}\setminus Q_K^1\subset\mathrm{lay}(K)$.
\end{definition}

\begin{lemma}[Fubini choice of the grid]\label{lem:fubini}
Let $\mathsf g:=|\nabla H_\alpha|^2\mathbf 1_D$. There is $\theta$ with
\[
\sum_{K\in\mathcal K_\theta}\int_{\mathrm{strip}^+(K)}\mathsf g\le\frac{24\tilde\ell}{R}C_An
\quad\text{and}\quad
\sum_{K\in\mathcal K_\theta}\int_{\mathrm{lay}(K)}\mathsf g\le\frac{24}{R}C_An .
\]
\end{lemma}
\begin{proof}
For $\theta$ uniform on $[0,R)^2$, a fixed $z$ lies in the strip of its square with probability
$\le|\{y\in[0,R)^2:\tilde\ell/2\le\operatorname{dist}(y,\partial[0,R]^2)\le2\tilde\ell+1\}|/R^2
\le4R(\tfrac32\tilde\ell+1)/R^2\le8\tilde\ell/R$, using $\tilde\ell\ge2$.
For the layer the probability is $\le8/R$. Take expectations, use Lemma~\ref{lem:apriori}(ii) and
apply Markov's inequality to each event with probability $\le1/3$.
\end{proof}

On $K\in\mathcal K_\theta$ we have $H_\alpha=H_{1/4}:=H-\sum_{\text{all }i}f_{1/4}(\cdot-z_i)$,
because points whose $\tfrac14$-disc meets $K$ have $d_i\ge\tfrac34$. Define
\[
\mathfrak F(Q_K):=\tfrac12\int_{Q_K}|\nabla H_{1/4}|^2+\sum_{z_i\in Q_K}\Bigl[\pi\log\tfrac14-2\pi\int_\C f_{1/4}(\cdot-z_i)m\Bigr].
\]

\begin{proposition}[Good squares]\label{prop:goodsquare}
Let $v_K$ be the density-freeze correction of Lemma~\ref{lem:freeze} on $Q_K$, let
$\delta_K:=\sup_{Q_K}|m-\bar m_K|$, and set
\[
S_K:=\int_{Q_K^{\tilde\ell/2}\setminus Q_K^{2\tilde\ell}}|\nabla H_{1/4}-\nabla v_K|^2 .
\]
Call $K$ \emph{good} if $S_K\le\ell^3\tilde\ell/C_{\rm scr}$, with $C_{\rm scr}$ as fixed in
Input~\ref{in:SCR} (above the ``Moreover'' constant of Prop.~7.18, so that $O$ is a concentric
square for every good $K$). Put
\[
c_*:=\Wtri+\tfrac\pi2\log4\pi+\pi\log\kappa .
\]
For every good $K$ and every $\epsilon\in(0,1]$,
\begin{multline*}
\mathfrak F(Q_K)+\pi\sum_{z_i\in Q_K}\log\lambda(x_i)\ge p_Kc_*
-C\Bigl[\frac{\ell S_K}{\tilde\ell}+R\tilde\ell+|\bar n_K-p_K|+p_K\delta_K\Bigr]\\
-\frac\epsilon2\int_{Q_K}|\nabla H_{1/4}|^2-\frac C\epsilon\int_{Q_K}|\nabla v_K|^2,
\end{multline*}
with $C$ depending only on $\kappa$.
\end{proposition}
\begin{proof}
\emph{Applying the screening input.} Apply Input~\ref{in:SCR} on $Q=Q_K$ with $\bar m=\bar m_K$
and $X=Z\cap Q_K$ (so $p=p_K$), and
\[
w:=H-\sum_{z_j\notin Q_K}f_{1/4}(\cdot-z_j)-v_K .
\]
The smeared outside charges are supported within $\tfrac14$ of $\partial Q_K$. So in $Q_K^{1/2}$,
$-\Delta w=2\pi(\sum\delta_{z_i}-m)-2\pi(\bar m_K-m)$, which is the required equation. Moreover
$w^\natural=H_{1/4}-v_K$ on $Q_K$, and $S\le S_K$. Since $R\ge R_0\ge10^5$,
$\max(4m_+/m_-,\ell_0)\le\ell\le\tilde\ell\le t_K/8$ (here $4m_+/m_-=100$, and $\ell\ge\ell_0$ because
$R\ge R_0\ge\ell_0^{5/2}$).
We obtain $Y_K$ with
\[
W^{\rm Neu}_{Q_K}(Y_K;\bar m_K)\le\mathcal B_K+C_S\bigl(\ell S_K/\tilde\ell+R\tilde\ell+|\bar n_K-p_K|\bigr),
\]
where $\mathcal B_K$ is the main term of Input~\ref{in:SCR}.

\emph{Comparing $\mathcal B_K$ with $\mathfrak F(Q_K)$.} Expand the square of
$\nabla H_{1/4}-\nabla v_K$ and use
$|\int\nabla H_{1/4}\cdot\nabla v_K|\le\frac\epsilon2\int|\nabla H_{1/4}|^2+\frac1{2\epsilon}\int|\nabla v_K|^2$.
For each $z_i\in Q_K$,
\[
\int_\C f_{1/4}(\cdot-z_i)m-\bar m_K\int_{Q_K}f_{1/4}(\cdot-z_i)\le\frac\pi{32}\Bigl(\delta_K+m_+\mathbf 1_{\operatorname{dist}(z_i,\partial Q_K)<1/4}\Bigr).
\]
By Lemma~\ref{lem:count}, at most $CR$ points are within $\tfrac14$ of $\partial Q_K$. Hence
\[
\mathcal B_K\le\mathfrak F(Q_K)+\frac\epsilon2\int_{Q_K}|\nabla H_{1/4}|^2+\Bigl(\frac1{2\epsilon}+\frac12\Bigr)\int|\nabla v_K|^2+Cp_K\delta_K+CR .
\]

\emph{Lower bound from Theorem~1.2 of the paper.} By Lemma~\ref{lem:reflection}(iv),
$W^{\rm Neu}\ge\bar n_K(\Wtri-\frac\pi2\log\bar m_K)$.

\emph{The chart factor.} $\pi\log\lambda(x_i)=\frac\pi2\log m(z_i)+\frac\pi2\log(4\pi\kappa^2)$, and
$|\log m(z_i)-\log\bar m_K|\le\delta_K/m_-$. Hence
\[
\pi\sum_{z_i\in Q_K}\log\lambda(x_i)\ge p_K\Bigl(\tfrac\pi2\log\bar m_K+\tfrac\pi2\log(4\pi\kappa^2)\Bigr)-Cp_K\delta_K .
\]

\emph{Conclusion.} Combine, using $|\log\bar m_K|\le C$ and $|\Wtri|<\infty$:
\[
\bar n_K\bigl(\Wtri-\tfrac\pi2\log\bar m_K\bigr)+p_K\tfrac\pi2\log\bar m_K\ge p_K\Wtri-C|\bar n_K-p_K| .
\]
\end{proof}

\begin{theorem}[Lower bound]\label{thm:LB}
For minimizers $Y$, $\mathfrak K(Y)\ge C_{\mathrm{tri}}\,n-o(n)$. That is,
$\liminf_n\min\mathfrak K/n\ge C_{\mathrm{tri}}$.
\end{theorem}
\begin{proof}
\emph{Setup.} Fix $R$ and take $n$ large. Choose $r$ by Lemma~\ref{lem:choose-r} and $\theta$ by
Lemma~\ref{lem:fubini}, and do the same in the north chart. We bound $\Phi_{\mathcal S}$; the
bound for $\Phi_{\mathcal N}$ is identical.

\emph{Partition of $D$.} By Lemma~\ref{lem:apriori}(i), and since $|\nabla H_\alpha|^2\ge0$ off
the good squares,
\[
\Phi_{\mathcal S}\ge\sum_{K\text{ good}}\Bigl[\mathfrak F(Q_K)+\pi\sum_{Q_K}\log\lambda\Bigr]
-C_1\#\{z_i\in D\setminus\textstyle\bigcup_{\rm good}Q_K\}-\pi\Lambda(r).
\]
The truncated form is additive over the disjoint sets $Q_K$ and their complement in $D$. On good
$Q_K$ we have $\alpha_i=\tfrac14$ and $H_\alpha=H_{1/4}$. Every other point contributes $\ge-C_1$,
as in Lemma~\ref{lem:apriori}(i).

\emph{Inputs to the error budget.} Write $\mathcal K=\mathcal K_\theta$; then
$\#\mathcal K\le|D|/R^2\le4\pi\kappa^2n/R^2$.
\begin{enumerate}[label=(\alph*)]
\item \emph{Strip energy.} By Lemma~\ref{lem:fubini} and Lemma~\ref{lem:freeze}, with
  $\delta_K\le C R/\sqrt n$ (since $|\nabla m|\le C/s_n$),
  \[
  \sum_KS_K\le2\sum_K\int_{\mathrm{strip}^+}\mathsf g+2\sum_K\int|\nabla v_K|^2
  \le C\frac{\tilde\ell}Rn+C\#\mathcal K\frac{R^6}{n}\le C\frac{\tilde\ell}Rn+CR^4 .
  \]
\item \emph{Bad squares.} Their number is $\le C_{\rm scr}\sum S_K/(\ell^3\tilde\ell)$. They
  contain at most $CR^2$ points each, so at most $C(R/\ell^3)n+C_R$ points in total.
\item \emph{Discrepancy.} By Lemma~\ref{lem:discrepancy} and Cauchy--Schwarz,
  \[
  \sum_K|\bar n_K-p_K|\le CR\,\#\mathcal K+\sqrt{2R\,\#\mathcal K\sum_K\int_{\mathrm{lay}(K)}\mathsf g}\le C\frac nR .
  \]
\item \emph{Uncovered points.} $D\setminus\bigcup_KQ_K$ consists of $D\setminus\bigcup_{\mathcal K}K$,
  which lies within $\sqrt2R+1$ of $\partial D$, and the sets $K\setminus Q_K$, each of area
  $\le2R$. By Lemma~\ref{lem:count} it contains $\le C(R\sqrt n+n/R)$ points.
\end{enumerate}

\emph{Summing over good squares.} Proposition~\ref{prop:goodsquare} with $\epsilon=R^2/\sqrt n$
gives
\begin{multline*}
\Phi_{\mathcal S}\ge c_*\,n_{\mathcal S}
-C\Bigl[\frac\ell{\tilde\ell}\Bigl(\frac{\tilde\ell}Rn+R^4\Bigr)+\frac{n}{R^2}R\tilde\ell+\frac nR+R\sqrt n\Bigr]
-C\frac{\epsilon}{2}C_An-\frac C\epsilon R^4\\
-(C_1+|c_*|)\Bigl[C\frac R{\ell^3}n+C_R+C\bigl(R\sqrt n+\tfrac nR\bigr)\Bigr]-\pi\Lambda(r).
\end{multline*}
Here $p_K\delta_K$ summed is $\le Cn\cdot R/\sqrt n$, and the bad and uncovered points pay
$c_*+C_1$ each.

\emph{Rates.} With $\ell=R^{2/5}$ and $\tilde\ell=R^{1/2}$ the coefficient of $n$ in the error is
$C(R^{-3/5}+R^{-1/2}+R^{-1}+R^{-1/5})$. The remaining terms are $O_R(\sqrt n)$. Hence
\[
\Phi_{\mathcal S}+\Phi_{\mathcal N}\ge c_*\,n-CR^{-1/5}n-C_R\sqrt n .
\]

\emph{Conclusion.} By Proposition~\ref{prop:twochart},
\[
\pi\mathfrak K\ge n\bigl(\Wtri+\tfrac\pi2\log4\pi\bigr)-CR^{-1/5}n-C_R\sqrt n .
\]
Divide by $\pi n$, let $n\to\infty$, then $R\to\infty$.
\end{proof}

\begin{remark}[How the technical points of the lower bound are handled]\label{rem:repairs}
\begin{itemize}[leftmargin=*]
\item \textbf{Screening error.} Screening is applied to the true potential (corrected by the freeze), so the
  error $\ell S_K/\tilde\ell$ is summed against Fubini-controlled strip energies. The parameters
  $\ell=R^{2/5}$, $\tilde\ell=R^{1/2}$ also make the bad-square fraction $R/\ell^3=R^{-1/5}$ vanish.
\item \textbf{The pole.} The pole is never in a chart's working region. The two-chart decomposition
  (Proposition~\ref{prop:twochart}) is exact, and every point is handled in a chart where
  $\lambda\in[\tfrac25,2]$.
\item \textbf{Quantization.} Quantization shrinks a Fubini-chosen grid square by at most one unit, and the
  inclusions of Definition~\ref{def:grid} transfer the strip control.
\item \textbf{Discrepancy sum.} Lemma~\ref{lem:discrepancy} together with step (c) of
  Theorem~\ref{thm:LB}.
\item \textbf{A priori bound.} Lemma~\ref{lem:apriori}(ii), which uses Theorem~\ref{thm:UB}.
\item \textbf{Filler interaction.} The filler interaction vanishes because $\eta=1$.
\item \textbf{Truncation comparison.} Separation at blown-up scale $\ge1$ makes all truncation
  radii equal $\tfrac14$, so the comparison of the energy with the truncated energy reduces to
  Lemma~\ref{lem:truncation} and Remark~\ref{rem:book-F}.
\end{itemize}
\end{remark}

\subsection{The filler on the screening layer}\label{sec:filler}

This subsection proves the filler bound used in Input~\ref{in:SCR}. Let $P\subset\C$ be compact
with piecewise $C^1$ boundary, $f\in L^\infty(P)$ and $g\in L^\infty(\partial P)$ (arclength). A
bounded vector field $E$ on $P$ \emph{solves $(f,g)$ weakly on $P$} if
\[
\int_PE\cdot\nabla\phi=\int_Pf\phi+\int_{\partial P}g\phi\dd\sigma\qquad\text{for every }\phi\in C^1(\bar P),
\tag{F1}
\]
that is, $-\operatorname{div}E=f$ in $P^\circ$ and $E\cdot\nu=g$ on $\partial P$. Taking $\phi\equiv1$
shows that $\int_Pf+\int_{\partial P}g=0$ is necessary.

\begin{lemma}[Explicit fields on a rectangle]\label{lem:rectfield}
Let $\mathrm{Rect}=a+[0,A]\times[0,B]$ with $b_-\le A,B\le b_+$. Let $f\in L^\infty(\mathrm{Rect})$ and
$g\in L^\infty(\partial\mathrm{Rect})$ with $\int_{\mathrm{Rect}}f+\int_{\partial\mathrm{Rect}}g=0$. Then there is a field
$E=(E_1,E_2)$ solving $(f,g)$ weakly on $\mathrm{Rect}$ with
\[
\|E\|_{L^\infty}\le3b_+\|f\|_{L^\infty}+14\,\frac{b_+}{b_-}\,\|g\|_{L^\infty}.
\]
In coordinates with $a=0$, $E_1(\cdot,y)$ is absolutely continuous on $[0,A]$ for every $y$, and
$E_2(x,\cdot)$ is absolutely continuous on $[0,B]$ for every $x$.
\end{lemma}
\begin{proof}
\emph{Construction.} Take $a=0$ and a \emph{bounded} Borel representative of $f$. Write
$g_W(y),g_E(y),g_S(x),g_N(x)$ for bounded Borel representatives of $g$ on the sides $\{x=0\}$,
$\{x=A\}$, $\{y=0\}$, $\{y=B\}$. Put
\[
k_x(x):=\int_0^Bf(x,y)\dd y+g_S(x)+g_N(x),\qquad k_y(y):=-g_W(y)-g_E(y),\qquad
G_{\rm tot}:=\int_0^Ak_x=\int_0^Bk_y ,
\]
where the last equality is the hypothesis $\int_{\mathrm{Rect}}f+\int_{\partial\mathrm{Rect}}g=0$. All these
functions are defined and bounded at \emph{every} $x$ and $y$. Let
$f_1(x,y):=k_x(x)/B+k_y(y)/A-G_{\rm tot}/(AB)$, so that $\int_0^Af_1(\cdot,y)=k_y(y)$ and
$\int_0^Bf_1(x,\cdot)=k_x(x)$. Set
\[
E_1(x,y):=-g_W(y)-\int_0^xf_1(\xi,y)\dd\xi,\qquad
E_2(x,y):=-g_S(x)-\int_0^y(f-f_1)(x,\eta)\dd\eta .
\]

\emph{Equation and fluxes.} $-\partial_xE_1=f_1$ and $-\partial_yE_2=f-f_1$ a.e. On the sides,
$-E_1(0,y)=g_W(y)$, $E_1(A,y)=-g_W(y)-k_y(y)=g_E(y)$, $-E_2(x,0)=g_S(x)$ and
$E_2(x,B)=-g_S(x)-\bigl(\int_0^Bf(x,\cdot)-k_x(x)\bigr)=g_N(x)$. For $\phi\in C^1(\bar{\mathrm{Rect}})$,
integrate $E_1\partial_x\phi$ along horizontal lines and $E_2\partial_y\phi$ along vertical lines
(fundamental theorem of calculus for absolutely continuous functions); this gives (F1).

\emph{Bounds.} Let $f_\infty=\|f\|_\infty$ and $g_\infty=\|g\|_\infty$. Then $|k_x|\le Bf_\infty+2g_\infty$,
$|k_y|\le2g_\infty$ and $|G_{\rm tot}|\le2Bg_\infty$, so $|f_1|\le f_\infty+2g_\infty/B+4g_\infty/A$. Hence
$|E_1|\le5g_\infty+Af_\infty+2(A/B)g_\infty$ and $|E_2|\le3g_\infty+2Bf_\infty+4(B/A)g_\infty$. Sum,
using $|E|\le|E_1|+|E_2|$.
\end{proof}

\paragraph{Point fields.} For $z\in\C$ and $r_1>0$ let
$\psi_z(x):=-\log(|x-z|/r_1)+(|x-z|^2-r_1^2)/(2r_1^2)$ for $0<|x-z|<r_1$ and $\psi_z(x):=0$ for
$|x-z|\ge r_1$. Then $\nabla\psi_z$ is continuous off $z$ and vanishes off $B(z,r_1)$. Since
$-\Delta\psi_z=2\pi(\delta_z-\mathbf 1_{B(z,r_1)}/(\pi r_1^2))$ and $\psi_z$, $\partial_r\psi_z$ vanish on
$\partial B(z,r_1)$,
\[
\int\nabla\psi_z\cdot\nabla\phi=2\pi\Bigl(\phi(z)-\frac1{\pi r_1^2}\int_{B(z,r_1)}\phi\Bigr)
\qquad(\phi\in C^1).
\tag{F2}
\]
For $0<\alpha\le r_1$, $\nabla f_\alpha(x-z)=-(x-z)/|x-z|^2$ on $B(z,\alpha)$. In polar coordinates the
integrand below is $r^2/r_1^4$ for $r<\alpha$ and $(1/r-r/r_1^2)^2$ for $\alpha<r<r_1$, so
\[
\int|\nabla\psi_z-\nabla f_\alpha(\cdot-z)|^2=2\pi\Bigl(\log\frac{r_1}\alpha+\frac{\alpha^2}{r_1^2}-\frac34\Bigr).
\tag{F3}
\]
Moreover $|\nabla\psi_z-\nabla f_\alpha(\cdot-z)|\le1/\alpha$ everywhere.

\paragraph{Variational energy $F^\flat$.} Let $U\subset\C$ be a bounded open set, $Z=(z_1,\dots,z_M)$ points
of $U$, $\tilde\mu\in L^\infty(U)$ with $\int_U\tilde\mu=M$, and let $\hat r_j$ be the radii (7.1.21) of $Z$ in
$U$. A \emph{weak competitor} for $(Z,\tilde\mu,U)$ is a field $E\in L^1_{\rm loc}(U;\R^2)$ such that
$E_{\hat r}:=E-\sum_j\nabla f_{\hat r_j}(\cdot-z_j)\in L^2(U)$ and
\[
\int_UE\cdot\nabla\phi=2\pi\Bigl(\sum_j\phi(z_j)-\int_U\phi\tilde\mu\Bigr)\qquad\text{for every }\phi\in C^1(\bar U).
\]
Define
\[
F^\flat(Z,\tilde\mu,U):=\inf_{E}\Bigl\{\frac1{4\pi}\int_U|E_{\hat r}|^2\Bigr\}-\frac12\sum_jg(\hat r_j)
-\sum_j\int_Uf_{\hat r_j}(x-z_j)\tilde\mu(x)\dd x ,
\]
with the infimum over weak competitors ($+\infty$ if there is none). This is the right-hand side of (7.1.22)
(with $c_d=2\pi$, as used in Lemma~\ref{lem:filler-energy}) with the Neumann gradient $\nabla u$ replaced
by a free weak competitor. If the Neumann solution $u$ of (7.1.1) exists on $U$, then $\nabla u$ is a weak
competitor and the projection Lemma~7.5 gives $\int_U|\nabla u_{\hat r}|^2\le\int_U|E_{\hat r}|^2$ for every
weak competitor $E$. Hence
\[
F(Z,\tilde\mu,U)\le F^\flat(Z,\tilde\mu,U).
\]
Defining $F^\flat$ requires no Neumann theory on $U$.

\stepcounter{theorem}
\begin{sublemma}[Filler field]\label{lem:filler}
Let $Q$ be a closed square of side $t$ and $O:=\{x\in Q:\operatorname{dist}(x,\partial Q)\ge w\}$, with
$w\ge\tfrac{10}3$ and $s:=t-2w\ge\tfrac{20}3$. Let $N_1:=\{x\in Q:\operatorname{dist}(x,O)\ge1\}$ (the
set $N_\eta$ of Prop.~7.18 with $\eta=1$ and $\Omega=Q$). Let $\tilde\mu\in L^\infty(N_1)$ with
$0\le\tilde\mu\le\tfrac1{100}$ a.e. and $M:=\int_{N_1}\tilde\mu\in\N$. Then there are $M$ points
$Z=(z_1,\dots,z_M)$ in the interior of $N_1$, with $|z_i-z_j|\ge\tfrac73$ for $i\ne j$ and
$\operatorname{dist}(z_i,\partial N_1)\ge\tfrac76$, and a field $E\in L^1(N_1;\R^2)$ such that
\begin{enumerate}[label=(\roman*)]
\item $\displaystyle\int_{N_1}E\cdot\nabla\phi=2\pi\Bigl(\sum_j\phi(z_j)-\int_{N_1}\phi\tilde\mu\Bigr)$ for every
  $\phi\in C^1(N_1)$;
\item $E_{1/4}:=E-\sum_j\nabla f_{1/4}(\cdot-z_j)\in L^\infty(N_1)$ and
  $\displaystyle\int_{N_1}|E_{1/4}|^2\le4\pi M\log(4r_1)+6\cdot10^7I$,
\end{enumerate}
where $r_1:=h_{\rm ring}/8$ and $I\le\tfrac9{70}|N_1|$ is the number of pieces, both defined in Step~1 below.
That is, $E$ is a weak competitor for $(Z,\tilde\mu,N_1^\circ)$, and every radius (7.1.21) of $Z$ in
$N_1^\circ$ equals $\tfrac14$.
\end{sublemma}
\begin{proof}
\emph{Step 1: rings and pieces.} Take $Q=[-\tfrac t2,\tfrac t2]^2$ and $O=[-\tfrac s2,\tfrac s2]^2$. Let
$\operatorname{dist}_\infty(x,O)=\max(|x_1|-\tfrac s2,|x_2|-\tfrac s2,0)$ be the sup-norm distance. Put
$n_r:=\lfloor3w/10\rfloor\ge1$, $h_{\rm ring}:=w/n_r\in[\tfrac{10}3,\tfrac{20}3)$ and $\xi_j:=jh_{\rm ring}$. The rings
$\mathrm{Ring}_j:=\{x\in N_1:\xi_{j-1}\le\operatorname{dist}_\infty(x,O)\le\xi_j\}$, $1\le j\le n_r$, cover $N_1$ and
overlap in null sets.
\begin{itemize}[leftmargin=*]
\item For $j\ge2$, $\mathrm{Ring}_j$ is the square annulus between the squares of half-side $\tfrac s2+\xi_{j-1}$
  and $\tfrac s2+\xi_j$. It consists of four corner squares of side $h_{\rm ring}$ and four straight strips of
  width $h_{\rm ring}$ and length $L_j:=s+2\xi_{j-1}$.
\item For $j=1$: next to a side of $O$, $\operatorname{dist}(x,O)=\operatorname{dist}_\infty(x,O)$, so the
  straight strips are, e.g., $\{\tfrac s2+1\le x_1\le\tfrac s2+h_{\rm ring},\ |x_2|\le\tfrac s2\}$, of width
  $h_{\rm ring}-1$ and length $L_1=s$. In the quadrant of a vertex $p$ of $O$, $\operatorname{dist}(x,O)=|x-p|$,
  so the corner parts are $K_p:=\mathrm{Rect}_p\setminus B(p,1)$, where $\mathrm{Rect}_p$ is the square of
  side $h_{\rm ring}$ with vertex $p$ in that quadrant (e.g.\ $\mathrm{Rect}_p=[\tfrac s2,\tfrac s2+h_{\rm ring}]^2$
  for $p=(\tfrac s2,\tfrac s2)$).
\end{itemize}
Cut each straight strip of $\mathrm{Ring}_j$ into $n_j:=\lfloor L_j/h_{\rm ring}\rfloor\ge1$ congruent rectangles
by cuts perpendicular to it; their length is $L_j/n_j\in[h_{\rm ring},2h_{\rm ring})$. The \emph{pieces} are
these rectangles, the corner squares of the rings $j\ge2$, and the four sets $K_p$. Every rectangular piece
has sides in $[h_{\rm ring}-1,2h_{\rm ring})$, hence aspect ratio $\le2h_{\rm ring}/(h_{\rm ring}-1)\le\tfrac{20}7$.
Every piece has area between $(h_{\rm ring}-1)h_{\rm ring}\ge\tfrac{70}9$ and $2h_{\rm ring}^2<89$ (also
$|K_p|=h_{\rm ring}^2-\pi/4$).

\emph{Order.} List the pieces as $P_1,\dots,P_I$, ring by ring from $j=1$. Within a ring go
counterclockwise from the lowest rectangle of the right strip: the right strip upwards, the
top-right corner piece, the top strip leftwards, the top-left corner piece, the left strip
downwards, the bottom-left corner piece, the bottom strip rightwards, the bottom-right corner piece.
Consecutive pieces $P_i,P_{i+1}$ share a segment $\sigma_i$, contained in a side of each, of length
$\ge h_{\rm ring}-1\ge\tfrac73$:
\begin{itemize}[leftmargin=*]
\item inside a strip, a full cross side; between a strip and a corner piece, the full end side of
  the strip's last rectangle;
\item for $K_p$, the parts at distance $\ge1$ from $p$ of the two sides of $\mathrm{Rect}_p$ through $p$; for
  the last corner piece of ring $1$, one such part and the opposite side of $\mathrm{Rect}_p$;
\item from the bottom-right corner piece of ring $j$ to the lowest rectangle of ring $j+1$,
  the segment $\{x_1=\tfrac s2+\xi_j,\ -\tfrac s2-\xi_j\le x_2\le-\tfrac s2-\xi_{j-1}\}$.
\end{itemize}
No $\sigma_i$ meets $B(p,1)$ for a vertex $p$ of $O$.

\emph{Step 2: masses, carries, points.} Let $m_i:=\int_{P_i}\tilde\mu\le|P_i|/100<\tfrac89$. Let
$M_0:=0$, $M_i:=m_1+\dots+m_i$, and
\[
k_i:=\lfloor M_i\rfloor-\lfloor M_{i-1}\rfloor,\qquad c_i:=M_i-\lfloor M_i\rfloor .
\]
Since $m_i<1$, $k_i\in\{0,1\}$. Moreover $c_i\in[0,1)$, $c_0=0$, $c_I=0$ (because $M_I=M\in\N$),
$\sum_ik_i=M$, and $c_i-c_{i-1}=m_i-k_i$. If $k_i=1$, put a point at the centre of $P_i$, or, for
$P_i=K_p$, at the centre of $\mathrm{Rect}_p$. Its distance to $\partial P_i$ is at least
$\min\bigl(\tfrac{h_{\rm ring}-1}2,\ \tfrac{h_{\rm ring}}{\sqrt2}-1\bigr)\ge\tfrac76$. Distinct points lie in
distinct pieces, so $|z_i-z_j|\ge\tfrac73$ and $\operatorname{dist}(z_i,\partial N_1)\ge\tfrac76$. In particular
every radius (7.1.21) of $Z$ in $U=N_1^\circ$ equals $\tfrac14$.

\emph{Step 3: local fields.} Let $r_1:=h_{\rm ring}/8\in[\tfrac5{12},\tfrac56)$, and
$\mathbb D_i:=B(z_i,r_1)\subset P_i^\circ$ when $k_i=1$. On $P_i$ put
\[
f_i:=2\pi\Bigl(k_i\frac{\mathbf 1_{\mathbb D_i}}{\pi r_1^2}-\tilde\mu\Bigr),\qquad
g_i:=\begin{cases}2\pi c_i/|\sigma_i|&\text{on }\sigma_i,\\-2\pi c_{i-1}/|\sigma_{i-1}|&\text{on }\sigma_{i-1},\\0&\text{on the rest of }\partial P_i.\end{cases}
\]
Then $\int_{P_i}f_i+\int_{\partial P_i}g_i=2\pi(k_i-m_i+c_i-c_{i-1})=0$,
$\|f_i\|_\infty\le2/r_1^2+2\pi/100<11.6$ and $\|g_i\|_\infty\le2\pi/(h_{\rm ring}-1)<2.7$.
\begin{itemize}[leftmargin=*]
\item \emph{Rectangular pieces.} Lemma~\ref{lem:rectfield}, with $b_+=2h_{\rm ring}<\tfrac{40}3$ and
  $b_+/b_-\le\tfrac{20}7$, gives $E_i^{\rm b}$ solving $(f_i,g_i)$ weakly on $P_i$ with
  $\|E_i^{\rm b}\|_\infty\le572$. So $\int_{P_i}|E_i^{\rm b}|^2\le2h_{\rm ring}^2\cdot572^2<3\cdot10^7$.
\item \emph{Corner pieces} $K=K_p=\mathrm{Rect}_p\setminus B(p,1)$. Let $q:=\mathrm{Rect}_p\cap B(p,1)$, extend
  $\tilde\mu$ by $0$ on $q$, and put
  $\tilde f:=2\pi(k_i\mathbf 1_{\mathbb D_i}/(\pi r_1^2)-\tilde\mu\mathbf 1_K)$ on $\mathrm{Rect}_p$. The datum
  $g_i$ vanishes on $\partial q\cap\partial\mathrm{Rect}_p$. Lemma~\ref{lem:rectfield} (with
  $b_+=b_-=h_{\rm ring}$, and bounded Borel representatives) gives $E^{\rm rect}$ solving $(\tilde f,g_i)$
  weakly on $\mathrm{Rect}_p$, with $\|E^{\rm rect}\|_\infty\le3h_{\rm ring}\cdot11.6+14\cdot2.7<270$. Write
  $x=p+re(\theta)$, $0\le\theta\le\pi/2$, in polar coordinates adapted to the quadrant of $\mathrm{Rect}_p$
  (e.g.\ $e(\theta)=(\cos\theta,\sin\theta)$ for $p=(\tfrac s2,\tfrac s2)$), with frame $e_r=e(\theta)$,
  $e_\theta=e'(\theta)$, and let
  $\gamma(\theta):=-E^{\rm rect}(p+e(\theta))\cdot e(\theta)$; this is defined for \emph{every} $\theta$
  (since $E^{\rm rect}$ is defined at every point), is bounded and Borel, and
  $|\gamma|<270$.

  The arc $\{|x-p|=1\}\cap\mathrm{Rect}_p$ is a graph over both axes. Integrating along horizontal and
  vertical lines as in the proof of Lemma~\ref{lem:rectfield}, now over $K$ and over $q$, gives for
  $\phi\in C^1(\bar{\mathrm{Rect}}_p)$
  \begin{align*}
  \int_KE^{\rm rect}\cdot\nabla\phi&=\int_K\tilde f\phi+\int_{\partial\mathrm{Rect}_p\cap\partial K}g_i\phi+\int_0^{\pi/2}\gamma(\theta)\phi(p+e(\theta))\dd\theta,\\
  \int_qE^{\rm rect}\cdot\nabla\phi&=-\int_0^{\pi/2}\gamma(\theta)\phi(p+e(\theta))\dd\theta .
  \end{align*}
  The second identity uses $\tilde f=0$ on $q$ and $g_i=0$ on $\partial q\cap\partial\mathrm{Rect}_p$. With
  $\phi\equiv1$ it gives $\int_0^{\pi/2}\gamma=0$.

  Let $\Upsilon(\theta):=\int_0^\theta\gamma$; it is Lipschitz with $\Upsilon(0)=\Upsilon(\pi/2)=0$. On the
  quarter annulus $A_q:=\{p+re(\theta):1\le r\le2,\ 0\le\theta\le\pi/2\}\subset K$ put
  \[
  G_{\rm cor}:=\gamma(\theta)\,\frac{2-r}r\,e_r+\Upsilon(\theta)\,e_\theta,\qquad G_{\rm cor}:=0\ \text{on }K\setminus A_q .
  \]
  In polar coordinates,
  $\int_KG_{\rm cor}\cdot\nabla\phi=\int\!\!\int\gamma(\theta)(2-r)\partial_r\phi\dd r\dd\theta+\int\!\!\int\Upsilon(\theta)\partial_\theta\phi\dd\theta\dd r
  =-\int_0^{\pi/2}\gamma(\theta)\phi(p+e(\theta))\dd\theta$: integrate by parts in $r$, using $2-r=0$ at
  $r=2$, and in $\theta$, using $\Upsilon'=\gamma$ and $\Upsilon(0)=\Upsilon(\pi/2)=0$. Hence
  $E_i^{\rm b}:=E^{\rm rect}|_K+G_{\rm cor}$ solves $(f_i,g_i)$ weakly on $K$, because the arc terms cancel. Moreover
  $|G_{\rm cor}|\le(1+\pi/2)\cdot270<700$, so
  $\int_K|E_i^{\rm b}|^2\le2h_{\rm ring}^2\cdot270^2+2\cdot\tfrac{3\pi}4\cdot700^2<10^7$.
\end{itemize}

\emph{Step 4: gluing.} Let $E:=\sum_i\mathbf 1_{P_i}\bigl(k_i\nabla\psi_{z_i}+E_i^{\rm b}\bigr)$, with
$\psi_{z_i}$ of radius $r_1$. It is in $L^1$, because $|\nabla\psi_z|\le1/|x-z|+C$. For $\phi\in C^1(N_1)$,
(F2) and the weak equations give
\begin{align*}
\int_{N_1}E\cdot\nabla\phi
&=\sum_i\Bigl[2\pi k_i\Bigl(\phi(z_i)-\frac{1}{\pi r_1^2}\int_{\mathbb D_i}\phi\Bigr)+\int_{P_i}f_i\phi+\int_{\partial P_i}g_i\phi\Bigr]\\
&=2\pi\Bigl(\sum_j\phi(z_j)-\int_{N_1}\tilde\mu\phi\Bigr).
\end{align*}
The averages over $\mathbb D_i$ cancel against the first part of $f_i$, and the boundary terms of $P_i$ and
$P_{i+1}$ on $\sigma_i$ are $+2\pi c_i/|\sigma_i|$ and $-2\pi c_i/|\sigma_i|$. So $E$ solves
$-\operatorname{div}E=2\pi(\sum_j\delta_{z_j}-\tilde\mu)$ in $N_1^\circ$ with $E\cdot\nu=0$ on $\partial N_1$,
weakly. This is (i). It is the gluing
argument of Corollary~7.6, (7.1.29) and (7.1.58), with
matched instead of vanishing normal fluxes on the interfaces $\sigma_i$.

\emph{Step 5: energy of $E_{1/4}$.} For $j\ne i$, $B(z_j,\tfrac14)\subset B(z_j,\tfrac76)\subset P_j$, so
on $P_i$ only $z_i$ (if $k_i=1$) contributes to $E_{1/4}$:
$E_{1/4}=k_i(\nabla\psi_{z_i}-\nabla f_{1/4}(\cdot-z_i))+E_i^{\rm b}$. By (F3) and the bound after it,
$E_{1/4}\in L^\infty$. Apply (F3) with $\alpha=\tfrac14\le r_1$; since $\alpha^2/r_1^2\le0.36<\tfrac34$,
\[
\int_{P_i}|E_{1/4}|^2\le4\pi k_i\log(4r_1)+2\int_{P_i}|E_i^{\rm b}|^2\le4\pi k_i\log(4r_1)+6\cdot10^7 .
\]
Sum over $i$ to get (ii).
\end{proof}

\begin{sublemma}[Filler energy]\label{lem:filler-energy}
In the situation of Lemma~\ref{lem:filler}, the points $Z$ satisfy
\[
F^\flat(Z,\tilde\mu,N_1^\circ)\le C_F|N_1|,\qquad C_F:=10^6 .
\]
\end{sublemma}
\begin{proof}
\emph{Competitor.} By Lemma~\ref{lem:filler}, $E$ is a weak competitor for $(Z,\tilde\mu,N_1^\circ)$, and
$E_{\hat r}=E_{1/4}$ because all radii (7.1.21) equal $\tfrac14$. Hence, with $c_d=2\pi$ and $g(\tfrac14)=\log4$
(as in (7.1.22)),
\[
F^\flat(Z,\tilde\mu,N_1^\circ)\le\frac1{4\pi}\int_{N_1}|E_{1/4}|^2-\frac M2\log4-\sum_j\int f_{1/4}(\cdot-z_j)\tilde\mu
\le M\log(2r_1)+\frac{6\cdot10^7}{4\pi}\,I ,
\]
where we used Lemma~\ref{lem:filler}(ii) and dropped the last sum, which is $\le0$ because $f_{1/4}\ge0$ and
$\tilde\mu\ge0$.

\emph{Numbers.} $\log(2r_1)=\log(h_{\rm ring}/4)<1$, $M\le|N_1|/100$ and $I\le\tfrac9{70}|N_1|$, so
$F^\flat(Z,\tilde\mu,N_1^\circ)\le(0.01+6.2\cdot10^5)|N_1|\le C_F|N_1|$.
\end{proof}

\paragraph{The energy $F$ of the Lectures.} By the inequality $F\le F^\flat$ of the paragraph ``Variational energy'',
Lemma~\ref{lem:filler-energy} also bounds the Neumann energy $F(Z,\tilde\mu,N_1^\circ)$, provided the
Neumann solution of (7.1.1) exists on $N_1^\circ$ (a non-convex domain with rounded corners). Nothing below
uses this: Input~\ref{in:SCR} uses $F^\flat$.

\begin{remark}[What the proof of Prop.~7.18 must provide]\label{rem:wp13-export}
For Input~\ref{in:SCR} (inner case, $d=2$, $U=\R^2$, $\Omega=Q$, $\eta=1$, $m:=m_-$) the proof of
Prop.~7.18 must provide, for the screened data of Input~\ref{in:SCR} (in particular for
$\ell\ge\ell_0(m_-,m_+)$, see Input~\ref{in:SCR}):
\begin{enumerate}[label=(X\arabic*)]
\item \emph{Shape of $O$} (``Moreover'' clause, with $C_{\rm scr}$ as in
  Input~\ref{in:SCR}, and $\ell\ge\ell_0$): $O=\{x\in Q:\operatorname{dist}(x,\partial Q)\ge w\}$ for some
  $w\in[\tilde\ell/2,2\tilde\ell]$.
\item \emph{The measure $\tilde\mu$} ((7.2.15)--(7.2.16), (A.0.15)--(A.0.16)): $\tilde\mu$ has a density on $N_\eta$, vanishes off $N_\eta$, satisfies
  $0\le\tilde\mu\le\bar m+m_-/2$ a.e.\ on $N_\eta$, and $\tilde\mu(N_\eta)=\bar n-n_O$.
\item[(X3$'$)] \emph{The bound (7.2.17)} for every configuration $Z$ of $\bar n-n_O$ points in $N_\eta$, with
  $J=\emptyset$ (Remark~7.19), where $F(Z,\tilde\mu,N_\eta)$ is replaced by the variational energy
  $F^\flat(Z,\tilde\mu,N_\eta^\circ)$ of \S\ref{sec:filler}.
\end{enumerate}
Items (X2) and (X3$'$) are the outputs of Appendix~A with the corrected equation
(Remark~\ref{rem:scr-interface}(d)).

\emph{Why $F^\flat$ suffices, and why (X4) is dropped.} By $F\le F^\flat$ (Lemma~7.5), (X3$'$) is implied by
(7.2.17) of the Lectures wherever the Neumann solution on $N_\eta^\circ$ exists. But Appendix~A uses
$F(Z,\tilde\mu,N_\eta)$ only through the field $E_3=\nabla h_3$, in (A.0.19) and (A.0.27),
and any field with the same weak divergence, zero normal flux and
the same energy serves there. So (X3$'$) is obtained from the proof of Appendix~A by taking for $E_3$ any
weak competitor, in place of the Neumann gradient. In this specialization the radii of $Z$ stay $\tfrac14$
($Z$ is $\tfrac73$-separated, at distance $\ge\tfrac{13}6$ from $O$ and $\ge\tfrac76$ from $\partial Q$), so
the enlargement step (7.1.40) is not needed for $E_3$. Consequently the former item (X4), the projection
Lemma~7.5 on $N_\eta^\circ$, and with it the Neumann problem on the non-convex domain $N_\eta^\circ$
(existence and regularity on a Lipschitz domain with rounded corners, and $H^1$-density of
$C^1(\bar U)$) are \emph{not} needed. The only projection lemma that remains in the proof of
Input~\ref{in:SCR} is the one on the square $Q$ (last bullet of Input~\ref{in:SCR}), where the Neumann
potential is explicit by reflection (Definition~\ref{def:neumann}, Lemma~\ref{lem:reflection}).

Not needed: the cells $H_k$ and constants $m_k$ of (A.0.14), the integrality of
$\int_{H_k^\eta}(\mu+m_k)$, a lower bound on $\tilde\mu$ beyond (X2), and Corollary~7.11. Given
(X1), (X2) and (X3$'$), Lemma~\ref{lem:filler-energy} applies: $\tilde\mu\le\tfrac32m_+\le\tfrac1{100}$,
$w\ge\tilde\ell/2\ge\tfrac{10}3$, and $s\ge t-4\tilde\ell\ge t/2\ge\tfrac{20}3$. The remaining dependence
on the screening proof at this step is therefore a matter of interface, not a gap.
\end{remark}

\begin{remark}[An alternative cutting of the filler layer]\label{rem:filler-route}
One could cut the cells $H_k^\eta$ of App.~A, which carry integer
$\tilde\mu$-mass by (A.0.14)--(A.0.15), into unit-mass rectangles with one centred point each, and glue
zero-flux competitors by (7.1.29). Away from the vertices of $O$ this works. But $N_\eta$ is
defined with the Euclidean distance (Prop.~7.18), so its inner boundary has four
quarter-circle arcs. The $H_k^\eta$ that meet them are not rectangles, and cutting them off
destroys the integrality of the cells. The proof above keeps the rest of that idea: pieces of
aspect ratio $\le\tfrac{20}7$, at most one point at the centre of each, an explicit competitor field
on each piece, and gluing as in (7.1.29). It cuts $N_\eta$ directly instead, and it moves the
fractional masses $c_i\in[0,1)$ between neighbouring pieces through their common sides. From
App.~A it uses only the density bound and the total mass.
\end{remark}

%======================================================================
\section{Inputs and dependencies}\label{sec:gaps}

\paragraph{External inputs.}
\begin{itemize}[leftmargin=*]
\item Theorem~1.2 of the paper (Input~\ref{in:R1}) and uniqueness of the torus potential (\lean{existsUnique\_isTorusPotential}).
\item (CF1), (CF2) (Input~\ref{in:CS}; proved in the closed-form note).
\item Screening (Input~\ref{in:SCR}; Prop.~7.18 and Rem.~7.19 of the Lectures, proved in the screening note), through the interface
  (X1), (X2), (X3$'$) of Remark~\ref{rem:wp13-export}. No Neumann theory on the non-convex layer $N_\eta^\circ$
  and no projection lemma on it is needed: the energy of the filler is the variational energy $F^\flat$,
  bounded by Lemma~\ref{lem:filler-energy} (the field is Lemma~\ref{lem:filler}). The only projection lemma
  used in Input~\ref{in:SCR} is the one on the square $Q$. Cor.~7.11 via (7.1.57) is no longer used.
\item Library facts: Neumann problem on squares and the Poincaré inequality
  (Lemma~\ref{lem:freeze}); the divergence theorem for Lipschitz functions on squares for a.e.
  level (Lemma~\ref{lem:discrepancy}); polar coordinates and circle averages
  (Lemma~\ref{lem:U}, Theorem~\ref{thm:separation}).
\item Further library facts:
  \begin{itemize}[leftmargin=*]
  \item Green's first identity on discs with finitely many punctures
    (Proposition~\ref{prop:splitting}, Lemma~\ref{lem:truncation});
  \item the circle-mean formula
    $\operatorname{avg}_{\partial B_r}\psi-\psi(x)=\int_0^r(2\pi s)^{-1}\int_{B_s}\Delta\psi\dd s$
    (Theorem~\ref{thm:separation});
  \item the circle average of $\log|z-a|$ (Lemma~\ref{lem:U}, Theorem~\ref{thm:separation});
  \item change of variables for the Dirichlet integral under conformal and anticonformal $C^2$
    maps (Lemma~\ref{lem:conformal});
  \item integrals of periodic functions agree over any fundamental domain, for the invariance of
    \texttt{IsTorusPotential} (Lemma~\ref{lem:reflection}(i));
  \item the removable singularity giving $u_Y-g(\cdot-y)\in C^1$ (admissibility in
    Definition~\ref{def:neumann}), from the theory of the torus potential.
  \end{itemize}
\item For Lemmas~\ref{lem:rectfield} and~\ref{lem:filler}: Fubini and the fundamental theorem of calculus for absolutely
  continuous functions on segments, and polar coordinates on a quarter annulus.
\end{itemize}

\paragraph{Points to check.}
\begin{itemize}[leftmargin=*]
\item The filler bound on thin layers in the screening input is proved by Lemmas~\ref{lem:filler} and~\ref{lem:filler-energy}. The remaining dependence is
  the interface (X1), (X2), (X3$'$) of Remark~\ref{rem:wp13-export}.
\item Interface checks for Input~\ref{in:SCR}: uniform constants, distinct interior
  output, the repaired proof, identification of the Neumann potential. These are statement-level
  matters, not mathematical gaps.
\item Constants such as $n_0$, $R_0$ and the $C$'s in Lemma~\ref{lem:geometry} and
  Lemma~\ref{lem:rowsum} are shown to exist but not computed. Each bound reduces to the explicit
  mass-coordinate formulas $\varrho(\tau(a))=\frac2na(n-a)$ and $|I_k|=\frac12\log(\cdots)$, with
  no further input. The polar-row cases (G4)--(G5) are the most delicate.
\end{itemize}
No other step depends on an unproved claim.

\paragraph{Errata and deviations.}
\begin{itemize}[leftmargin=*]
\item The misprint in (5.3.4) of the Lectures is confirmed by an independent reading (Remark~\ref{rem:E2}).
\item The error in (7.1.37) of the Lectures affects only the screening input (Remark~\ref{rem:scr-interface}(d)). The error in Appendix~A
  ($-\Delta h_{2,k}$ for $\eta>0$) also affects only the screening input: (X2)/(X3$'$) come from Appendix~A with the corrected equation
  (Remark~\ref{rem:scr-interface}(d)). Both are listed in the errata record of the repository.
\item Prop.~7.18 as printed omits the hypothesis $\ell\ge\ell_0(m_-,m_+)$ used by Appendix~A; it is added
  to Input~\ref{in:SCR}.
\item The filler bound deviates from the route of the Lectures, which bound the filler energy by
  (7.1.57); here Lemmas~\ref{lem:filler} and~\ref{lem:filler-energy} prove it directly
  (Remark~\ref{rem:filler-route}).
\item A construction by mirrored hemispheres produces only even $n$. Construction~\ref{con:rings}
  produces every $n$.
\item Stationarity (Remark~\ref{rem:stationarity}) is not used.
\end{itemize}

\begin{thebibliography}{9}
\bibitem{Preprint} OpenAI, \emph{Triangular minimality for planar Coulomb renormalized energy}, preprint, 23~September 2026.
\bibitem{Lectures} S.~Serfaty, \emph{Lectures on Coulomb and Riesz gases}, arXiv:2407.21194, 2024.
\bibitem{BeterminSandier2018} L.~B\'etermin and \'E.~Sandier, \emph{Renormalized energy and asymptotic expansion of optimal
logarithmic energy on the sphere}, Constr.\ Approx.\ \textbf{47} (2018), 39--74.
\bibitem{BrauchartHardinSaff2012} J.~S.~Brauchart, D.~P.~Hardin and E.~B.~Saff, \emph{The next-order term for optimal Riesz and
logarithmic energy asymptotics on the sphere}, in: Recent Advances in Orthogonal Polynomials, Special Functions, and Their
Applications, Contemp.\ Math.\ \textbf{578}, Amer.\ Math.\ Soc., 2012, 31--61.
\end{thebibliography}

\end{document}
